\documentclass[12pt,a4paper]{article}

% ── Encoding & Sprache ────────────────────────────────────────────────────────
\usepackage[T1]{fontenc}
\usepackage[utf8]{inputenc}
\usepackage[ngerman]{babel}

% ── Typographie ───────────────────────────────────────────────────────────────

\usepackage{rotating}
\usepackage{microtype}
\usepackage{setspace}
\onehalfspacing

% ── Seitenränder ──────────────────────────────────────────────────────────────
\usepackage[left=3cm,right=2.5cm,top=2.5cm,bottom=2.5cm]{geometry}

% ── Mathematik ────────────────────────────────────────────────────────────────
\usepackage{amsmath,amssymb,amsthm}
\usepackage{stmaryrd}

% ── Grafik & Farbe ────────────────────────────────────────────────────────────
\usepackage{graphicx}
\usepackage[table,xcdraw]{xcolor}
\usepackage{tikz}
\usetikzlibrary{shapes.geometric,arrows.meta,positioning,fit,backgrounds,calc}

% ── Tabellen ──────────────────────────────────────────────────────────────────
\usepackage{booktabs}
\usepackage{longtable}
\usepackage{multirow}
\usepackage{array}
\usepackage{tabularx}

% ── Code-Listing ──────────────────────────────────────────────────────────────
\usepackage{listings}
\usepackage{listingsutf8}

% Farben für Listings
\definecolor{codegreen}{rgb}{0.0,0.55,0.0}
\definecolor{codegray}{rgb}{0.5,0.5,0.5}
\definecolor{codepurple}{rgb}{0.55,0.0,0.55}
\definecolor{codeblue}{rgb}{0.0,0.0,0.75}
\definecolor{codeback}{rgb}{0.97,0.97,0.97}
\definecolor{framecol}{rgb}{0.7,0.7,0.7}
\definecolor{sqlback}{rgb}{0.95,0.98,1.0}
\definecolor{pyback}{rgb}{0.97,0.99,0.97}
\definecolor{javaback}{rgb}{1.0,0.97,0.95}
\definecolor{csback}{rgb}{0.97,0.95,1.0}

% Basis-Style
\lstdefinestyle{basestyle}{
  basicstyle=\ttfamily\footnotesize,
  keywordstyle=\color{codeblue}\bfseries,
  commentstyle=\color{codegreen}\itshape,
  stringstyle=\color{codepurple},
  numberstyle=\tiny\color{codegray},
  numbers=left,
  stepnumber=1,
  numbersep=8pt,
  tabsize=2,
  showspaces=false,
  showstringspaces=false,
  breaklines=true,
  breakatwhitespace=true,
  frame=single,
  rulecolor=\color{framecol},
  captionpos=b,
  aboveskip=8pt,
  belowskip=8pt,
}

% SQL
\lstdefinestyle{sqlstyle}{
  style=basestyle,
  backgroundcolor=\color{sqlback},
  language=SQL,
  morekeywords={SELECT,FROM,WHERE,JOIN,LEFT,RIGHT,INNER,OUTER,
    GROUP,BY,ORDER,HAVING,AS,ON,LIMIT,OFFSET,DISTINCT,IN,
    NOT,AND,OR,EXISTS,BETWEEN,LIKE,IS,NULL,COUNT,SUM,AVG,
    MIN,MAX,OVER,PARTITION,ROWS,RANGE,UNBOUNDED,PRECEDING,
    FOLLOWING,CURRENT,ROW,CASE,WHEN,THEN,ELSE,END,WITH,
    UNION,INTERSECT,EXCEPT,ALL,ANY,SOME,INSERT,UPDATE,DELETE,
    CREATE,TABLE,INDEX,VIEW,TRIGGER,CASCADE,CONSTRAINT},
  keywordstyle=\color{codeblue}\bfseries,
}

% Python
\lstdefinestyle{pythonstyle}{
  style=basestyle,
  backgroundcolor=\color{pyback},
  language=Python,
  morekeywords={lambda,yield,from,import,as,with,pass,
    True,False,None,self,cls},
}

% Java
\lstdefinestyle{javastyle}{
  style=basestyle,
  backgroundcolor=\color{javaback},
  language=Java,
  morekeywords={var,record,sealed,permits,yield,instanceof,
    stream,filter,map,flatMap,collect,toList,sorted,distinct,
    limit,skip,reduce,count,anyMatch,allMatch,noneMatch,
    findFirst,findAny,min,max,sum,average,groupingBy,
    partitioningBy,joining,toMap,Collectors,Optional},
}

% C# / LINQ
\lstdefinelanguage{CSharp}[]{Java}{
  morekeywords={var,from,where,select,group,by,join,on,equals,
    into,orderby,ascending,descending,let,in,yield,return,
    async,await,record,init,with,new,namespace,using,class,
    interface,struct,enum,void,string,int,bool,double,decimal,
    List,IEnumerable,IQueryable,Task,Func,Action,Predicate},
}
\lstdefinestyle{csharpstyle}{
  style=basestyle,
  backgroundcolor=\color{csback},
  language=CSharp,
}

% ── Hyperlinks & PDF-Metadaten ────────────────────────────────────────────────
\usepackage[
  colorlinks=true,
  linkcolor=codeblue,
  citecolor=codeblue,
  urlcolor=codeblue,
  pdftitle={UQTL -- Unified Query Translation Layer},
  pdfauthor={Stephan},
  pdfsubject={Standardisierung der sprachübergreifenden Query-Übersetzung},
]{hyperref}

% ── Inhaltsverzeichnis ────────────────────────────────────────────────────────
\usepackage{tocloft}
\setlength{\cftbeforesecskip}{4pt}

% ── Theoreme & Definitionen ───────────────────────────────────────────────────
\theoremstyle{definition}
\newtheorem{definition}{Definition}[section]
\newtheorem{beispiel}{Beispiel}[section]
\newtheorem{regel}{Übersetzungsregel}[section]
\theoremstyle{plain}
\newtheorem{theorem}{Theorem}[section]
\newtheorem{lemma}[theorem]{Lemma}
\theoremstyle{remark}
\newtheorem{bemerkung}{Bemerkung}[section]

% ── Abkürzungen ───────────────────────────────────────────────────────────────
\usepackage{acronym}

% ── Caption-Formatierung ──────────────────────────────────────────────────────
\usepackage[labelfont=bf,textfont=it]{caption}
\usepackage{subcaption}

% ── Sonstige Hilfspakete ──────────────────────────────────────────────────────
\usepackage{enumitem}
\usepackage{parskip}
\usepackage{mdframed}
\usepackage{float}
\usepackage{pdflscape}
\usepackage{afterpage}

% Hervorhebungsbox
\newmdenv[
  backgroundcolor=blue!5,
  linecolor=codeblue,
  linewidth=1pt,
  roundcorner=4pt,
  innerleftmargin=10pt,
  innerrightmargin=10pt,
  innertopmargin=8pt,
  innerbottommargin=8pt,
]{infobox}

\newmdenv[
  backgroundcolor=green!5,
  linecolor=codegreen,
  linewidth=1pt,
  roundcorner=4pt,
  innerleftmargin=10pt,
  innerrightmargin=10pt,
  innertopmargin=8pt,
  innerbottommargin=8pt,
]{successbox}

\newmdenv[
  backgroundcolor=orange!5,
  linecolor=orange!70,
  linewidth=1pt,
  roundcorner=4pt,
  innerleftmargin=10pt,
  innerrightmargin=10pt,
  innertopmargin=8pt,
  innerbottommargin=8pt,
]{warningbox}

% ── Titelseite ────────────────────────────────────────────────────────────────
\title{%
  \vspace{-1cm}%
  {\large Technischer Standard -- Entwurf v2.0}\\[0.5em]
  \rule{\linewidth}{0.6pt}\\[0.5em]
  {\LARGE\bfseries UQTL -- Unified Query Translation Layer}\\[0.3em]
  {\large Sprachübergreifende Standardisierung der Query"=Übersetzung\\
  von Python, Java (Streams) und C\#/.NET (LINQ) nach SQL}\\[0.5em]
  \rule{\linewidth}{0.6pt}%
}
\author{%
  Stephan Epp\\[0.3em]
  \small Datenbankstandards, Sprachdesign, API-Normierung
}
\date{6. März 2026}

% ═══════════════════════════════════════════════════════════════════════════════
\begin{document}
% ═══════════════════════════════════════════════════════════════════════════════

\maketitle
\thispagestyle{empty}

% ── Abstract ──────────────────────────────────────────────────────────────────
\begin{abstract}
\noindent
Die vorliegende Arbeit definiert den \textbf{Unified Query Translation Layer (UQTL)},
einen offenen technischen Standard zur einheitlichen Übersetzung von
datenabfragenden Ausdrücken aus den Programmiersprachen Python, Java (Stream API)
und C\#/.NET (LINQ) in äquivalente SQL-Anweisungen. Das zentrale Ziel des Standards
ist die \emph{semantische Äquivalenzgarantie}: Gleichartige Abfragen, unabhängig
von der Quellsprache formuliert, erzeugen stets dasselbe normalisierte SQL-Statement.

Der Standard umfasst eine formale Grammatik für den sprachunabhängigen intermediären
Ausdrucksbaum (Intermediate Query Tree, IQT), vollständige Übersetzungsregeln
für alle gängigen Query-Operatoren (Filterung, Projektion, Aggregation,
Sortierung, Joins, Fensterfunktionen und Unterabfragen), sowie Konformitätsstufen
(Conformance Levels CL-1 bis CL-4) für Implementierungen. Umfangreiche
Vergleichsbeispiele belegen die Anwendbarkeit des Standards auf reale
Datenbankszenarien.

\medskip
\noindent\textbf{Schlagwörter:} Query Translation, SQL-Generierung, LINQ, Python ORM,
Java Streams, Datenbankabstraktion, AST-Normierung, Interoperabilität
\end{abstract}

\newpage
\tableofcontents
\newpage

% ─────────────────────────────────────────────────────────────────────────────
% ═══════════════════════════════════════════════════════════════════════════════
\section{Einleitung}
\label{sec:einleitung}
% ═══════════════════════════════════════════════════════════════════════════════

\subsection{Motivation und Problemstellung}
\label{subsec:motivation}

Moderne Softwaresysteme sind in hohem Maße polyglott: Unterschiedliche Dienste
einer verteilten Anwendungslandschaft werden in verschiedenen Programmiersprachen
entwickelt, greifen jedoch häufig auf gemeinsame relationale Datenbanksysteme zu.
Python-basierte Datenpipelines, Java-Microservices und C\#-Backend-Dienste
formulieren dabei konzeptuell identische Datenbankabfragen -- verwenden jedoch
vollständig unterschiedliche Ausdrucksmittel.

Betrachten wir die simple Frage \emph{,,Alle aktiven Nutzer, deren Umsatz im
vergangenen Quartal über 1\,000~Euro lag, sortiert nach Name''}: In Python
(SQLAlchemy), Java (Stream API mit JPA Criteria) und C\#/.NET (LINQ to Entities)
entstehen syntaktisch völlig unterschiedliche Ausdrücke, die jedoch -- bei
korrekter Implementierung -- denselben SQL-Ausdruck erzeugen sollten. In der
Praxis tun sie dies \emph{nicht}: ORMs und Query-Builder unterschiedlicher
Sprachen und Frameworks produzieren inkompatible, unvollständige oder semantisch
abweichende SQL-Ausgaben \cite{ansi_sql_2016,hibernate_docs,ef_core_docs}.

Daraus entstehen konkrete Probleme in der Softwareentwicklung:

\begin{itemize}[leftmargin=1.5cm]
  \item \textbf{Inkonsistente Abfrageergebnisse:} Zwei Dienste, die dieselbe
    Geschäftslogik ausdrücken, erhalten unterschiedliche Daten.
  \item \textbf{Erschwerte Code-Reviews:} Entwickler müssen den mentalen Overhead
    betreiben, Query-Ausdrücke verschiedener Sprachen zu verstehen und zu vergleichen.
  \item \textbf{Migrations-Risiken:} Bei der Portierung von Java- nach Python-Diensten
    oder umgekehrt kann keine mechanische Äquivalenzprüfung erfolgen.
  \item \textbf{Fehlende Testbarkeit:} Ohne definiertes Normalformat ist die
    automatisierte Verifikation der Query-Äquivalenz nicht möglich.
\end{itemize}

\subsection{Zielsetzung des Standards}
\label{subsec:ziel}

Der \textbf{Unified Query Translation Layer (UQTL)} adressiert diese Problematik
durch die Definition eines umfassenden technischen Standards, der folgende
Ziele verfolgt:

\begin{enumerate}[leftmargin=1.5cm]
  \item \textbf{Semantische Äquivalenzgarantie:} Ausdrücke derselben logischen
    Abfrage in Python, Java und C\# werden stets in dasselbe normalisierte
    SQL-Statement übersetzt.
  \item \textbf{Intermediäre Darstellung:} Ein sprachunabhängiger
    \emph{Intermediate Query Tree (IQT)} dient als gemeinsame Zwischenrepräsentation.
  \item \textbf{Vollständige Operatorabdeckung:} Alle wesentlichen
    Query-Operatoren (Filter, Projektion, Aggregation, Sortierung, Joins,
    Fensterfunktionen, Unterabfragen) sind standardisiert.
  \item \textbf{Erweiterbarkeit:} Der Standard definiert Erweiterungspunkte
    für datenbankspezifische Dialekte (PostgreSQL, MySQL, Oracle, MSSQL).
  \item \textbf{Konformitätsstufen:} Vier Konformitätsstufen (CL-1 bis CL-4)
    erlauben schrittweise Implementierungen.
\end{enumerate}

\begin{infobox}
\textbf{Kerndefinition des Standards:}

Für jede Menge logisch äquivalenter Query-Ausdrücke $Q_{py}$, $Q_{java}$,
$Q_{cs}$ aus Python, Java bzw.\ C\# gilt unter UQTL:
\[
  \mathrm{SQL}(Q_{py}) \;\equiv_\sigma\; \mathrm{SQL}(Q_{java})
  \;\equiv_\sigma\; \mathrm{SQL}(Q_{cs})
\]
wobei $\equiv_\sigma$ die semantische SQL-Äquivalenz unter einer gegebenen
Datenbankschema-Instanz $\sigma$ bezeichnet.
\end{infobox}

\subsection{Abgrenzung}
\label{subsec:abgrenzung}

Der UQTL-Standard ist \emph{kein} vollständiges ORM-Framework und definiert
\emph{keine} Laufzeitbibliothek. Er spezifiziert:

\begin{itemize}[leftmargin=1.5cm]
  \item Die formale Syntax und Semantik des Intermediate Query Tree (IQT)
  \item Ãœbersetzungsregeln von den drei Quellsprachen in den IQT
  \item Ãœbersetzungsregeln vom IQT nach ANSI SQL:2016
  \item Normalisierungsregeln für eindeutige SQL-Ausgaben
  \item Konformitätsanforderungen für Implementierungen
\end{itemize}

Explizit nicht Teil des Standards sind: Datenbankzugriff, Verbindungsverwaltung,
Entity-Mapping, Caching-Strategien sowie dialektspezifische Optimierungen
(diese können als optionale Erweiterungen definiert werden).

\subsection{Aufbau der Arbeit}
\label{subsec:aufbau}

\ref{sec:grundlagen} legt die theoretischen Grundlagen dar:
relationale Algebra, Query-Abstraktionen und bestehende Ansätze.
\ref{sec:iqt} definiert den Intermediate Query Tree formal.
\ref{sec:operatoren} spezifiziert alle Ãœbersetzungsregeln.
\ref{sec:beispiele} demonstriert den Standard anhand umfangreicher
Praxisbeispiele.
\ref{sec:konformitaet} definiert die Konformitätsstufen.
\ref{sec:implementierung} beschreibt Referenzimplementierungen.
\ref{sec:diskussion} diskutiert Limitierungen und Erweiterungen.
\ref{sec:zusammenfassung} fasst die Ergebnisse zusammen.


\newpage
% ═══════════════════════════════════════════════════════════════════════════════
\section{Theoretische Grundlagen}
\label{sec:grundlagen}
% ═══════════════════════════════════════════════════════════════════════════════

\subsection{Relationale Algebra und SQL}
\label{subsec:rel_algebra}

Die theoretische Basis für die Query-Übersetzung bildet die \emph{relationale
Algebra} nach Codd \cite{codd1970relational}. Alle im UQTL-Standard behandelten
Query-Operatoren lassen sich auf eine endliche Menge algebraischer
Grundoperationen zurückführen.

\begin{definition}[Relationale Grundoperationen]
\label{def:rel_ops}
Gegeben seien Relationen $R$, $S$ über einem Schema $\mathcal{R}$.
Die relationalen Grundoperationen sind:
\begin{align*}
  \sigma_p(R) &\quad \text{Selektion (Filterung) mit Prädikat } p \\
  \pi_A(R)    &\quad \text{Projektion auf Attributmenge } A \subseteq \mathcal{R} \\
  R \times S  &\quad \text{Kartesisches Produkt} \\
  R \cup S    &\quad \text{Vereinigung} \\
  R - S       &\quad \text{Differenz} \\
  R \bowtie_p S &\quad \text{Join mit Prädikat } p \\
  \gamma_{A,F}(R) &\quad \text{Gruppierung nach } A \text{ mit Aggregatfunktionen } F
\end{align*}
\end{definition}

SQL ist eine deklarative Anfragesprache, die über diese algebraischen
Grundoperationen hinausgeht und um Fensterfunktionen, rekursive Abfragen
(CTE), Mengenoperationen sowie datenbankspezifische Erweiterungen ergänzt wird.
ANSI SQL:2016 \cite{ansi_sql_2016} bildet die Ziel-Spezifikation des UQTL-Standards.

\subsection{Query-Abstraktion in modernen Programmiersprachen}
\label{subsec:query_abstraktion}

Alle drei im UQTL-Standard behandelten Sprachen implementieren das Konzept der
\emph{Query Comprehension} oder \emph{monadischen Query-Komposition}:

\subsubsection{Python -- SQLAlchemy Core \& ORM}

Python verfügt über keine native Query-Syntax in der Sprache selbst.
Der De-facto-Standard für relationale Abfragen ist \textbf{SQLAlchemy},
dessen \emph{Core}-Schicht einen Expression Language API bereitstellt.
Abfragen werden als verkettete Methodenaufrufe oder als ORM-Ausdrücke
formuliert \cite{sqlalchemy_docs}:

\begin{lstlisting}[style=pythonstyle, caption={Python SQLAlchemy -- Query-Stilarten},
  label=lst:py_styles]
# Stil 1: SQLAlchemy ORM (Session.query -- Legacy)
users = session.query(User)\
    .filter(User.active == True)\
    .order_by(User.name)\
    .all()

# Stil 2: SQLAlchemy 2.0 Select-Style (modern)
stmt = select(User)\
    .where(User.active == True)\
    .order_by(User.name)
users = session.execute(stmt).scalars().all()

# Stil 3: Pandas / pandasql (fuer Datenpipelines)
import pandasql as ps
users = ps.sqldf("SELECT * FROM users WHERE active=1 ORDER BY name")
\end{lstlisting}

\subsubsection{Java -- Stream API und JPA Criteria API}

Java bietet seit Version 8 die \textbf{Stream API} für funktionale
Datenverarbeitung auf Collections \cite{java_stream_docs}. Für
Datenbankabfragen wird die \textbf{JPA Criteria API} oder Frameworks
wie \textbf{QueryDSL} verwendet:

\begin{lstlisting}[style=javastyle, caption={Java -- Query-Stilarten},
  label=lst:java_styles]
// Stil 1: Stream API (In-Memory-Verarbeitung)
List<User> users = allUsers.stream()
    .filter(u -> u.isActive())
    .sorted(Comparator.comparing(User::getName))
    .collect(Collectors.toList());

// Stil 2: JPA Criteria API (Datenbankabfrage)
CriteriaBuilder cb = em.getCriteriaBuilder();
CriteriaQuery<User> cq = cb.createQuery(User.class);
Root<User> root = cq.from(User.class);
cq.select(root)
  .where(cb.isTrue(root.get("active")))
  .orderBy(cb.asc(root.get("name")));

// Stil 3: QueryDSL (typsicher, fliessend)
List<User> users = new JPAQuery<>(em)
    .select(QUser.user)
    .from(QUser.user)
    .where(QUser.user.active.isTrue())
    .orderBy(QUser.user.name.asc())
    .fetch();
\end{lstlisting}

\subsubsection{C\# -- LINQ (Language Integrated Query)}

C\# bietet mit \textbf{LINQ} eine nativ in die Sprache integrierte
Query-Syntax, die vom Compiler direkt in Methodenaufrufe übersetzt wird.
Entity Framework Core implementiert den \texttt{IQueryable<T>}-Provider
für SQL-Generierung \cite{ef_core_docs,linq_docs}:

\begin{lstlisting}[style=csharpstyle, caption={C\# LINQ -- Query-Stilarten},
  label=lst:cs_styles]
// Stil 1: LINQ Query Syntax (deklarativ, SQL-aehnlich)
var users = from u in context.Users
            where u.Active == true
            orderby u.Name
            select u;

// Stil 2: LINQ Method Syntax (fliessend, Lambda)
var users = context.Users
    .Where(u => u.Active == true)
    .OrderBy(u => u.Name)
    .ToList();

// Stil 3: LINQ mit anonymen Typen (Projektion)
var result = context.Users
    .Where(u => u.Active)
    .Select(u => new { u.Id, u.Name, u.Email });
\end{lstlisting}

\subsection{Bestehende Ansätze und deren Limitierungen}
\label{subsec:bestehend}

Tabelle~\ref{tab:bestehend} fasst bestehende Ansätze zur Query-Abstraktion
und deren Grenzen zusammen.

\begin{table}[H]
\centering
\caption{Bestehende Query-Abstraktionsansätze und ihre Limitierungen}
\label{tab:bestehend}
\renewcommand{\arraystretch}{1.3}
\begin{tabularx}{\textwidth}{lXX}
\toprule
\textbf{Ansatz} & \textbf{Merkmale} & \textbf{Limitierungen} \\
\midrule
Hibernate (Java)
  & Mächtiges ORM, HQL, Criteria API
  & Nur Java; erzeugt oft ineffiziente SQL; kein sprachübergreifender Standard \\
\hline
Entity Framework (C\#)
  & LINQ-Integration, Migrations
  & Nur .NET; SQL-Ausgabe variiert je nach EF-Version \\
\hline
SQLAlchemy (Python)
  & Flexibel, Core + ORM
  & Nur Python; verschiedene API-Generationen; inkonsistente SQL-Ausgabe \\
\hline
JOOQ (Java)
  & Typsicherer SQL-Builder
  & Nur Java; erfordert Code-Generierung; kommerziell für Enterprise \\
\hline
Exposed (Kotlin)
  & Kotlin DSL für SQL
  & Nur JVM/Kotlin; junges Ökosystem \\
\hline
Prisma (Node.js)
  & Typsicherer ORM für TypeScript
  & Nur Node.js/TypeScript; anderes Ökosystem \\
\hline
GraphQL-to-SQL
  & Graphenabfragen über SQL
  & Anderes Paradigma; nicht auf Standard-CRUD ausgerichtet \\
\bottomrule
\end{tabularx}
\end{table}

\begin{warningbox}
\textbf{Kritische Feststellung:}
Kein bestehender Ansatz garantiert sprachübergreifende SQL-Äquivalenz.
Die Lücke liegt nicht in der fehlenden Implementierung, sondern im Fehlen
eines formalen Standards, der Äquivalenzanforderungen und
Normalisierungsregeln verbindlich definiert.
\end{warningbox}

\subsection{Formale Grundbegriffe}
\label{subsec:formale_grundbegriffe}

\begin{definition}[Query-Ausdruck]
\label{def:query_expr}
Ein \emph{Query-Ausdruck} $Q$ ist ein Term einer Operatoralgebra $\mathcal{A}$
über einer Datenbankschema-Instanz $\sigma$. Die \emph{Semantik} von $Q$ ist
die Relation $\llbracket Q \rrbracket_\sigma$, die durch Auswertung von $Q$
in $\sigma$ entsteht.
\end{definition}

\begin{definition}[Semantische Äquivalenz]
\label{def:sem_aequiv}
Zwei Query-Ausdrücke $Q_1$ und $Q_2$ heißen \emph{semantisch äquivalent},
geschrieben $Q_1 \equiv Q_2$, wenn für alle Datenbankinstanzen $\sigma$ gilt:
\[
  \llbracket Q_1 \rrbracket_\sigma = \llbracket Q_2 \rrbracket_\sigma
\]
\end{definition}

\begin{definition}[UQTL-Konformität]
\label{def:konformitaet}
Ein Übersetzungssystem $T$ heißt \emph{UQTL-konform}, wenn für alle
äquivalenten Quellanfragen $Q_{py} \equiv Q_{java} \equiv Q_{cs}$ gilt:
\[
  T(Q_{py}) = T(Q_{java}) = T(Q_{cs}) = S^*
\]
wobei $S^*$ das normalisierte SQL-Statement im UQTL-Normalformat ist.
\end{definition}

Diese Definitionen bilden das formale Fundament des gesamten Standards.


\newpage
% ═══════════════════════════════════════════════════════════════════════════════
\section{Der Intermediate Query Tree (IQT)}
\label{sec:iqt}
% ═══════════════════════════════════════════════════════════════════════════════

\subsection{Überblick und Architektur}
\label{subsec:iqt_uebersicht}

Das Herzstück des UQTL-Standards ist der \textbf{Intermediate Query Tree
(IQT)}: eine sprachunabhängige Baumstruktur, die jeden Query-Ausdruck aus
Python, Java und C\# einheitlich repräsentiert. Der IQT fungiert als
\emph{Pivot-Sprache} in der Übersetzungspipeline.

\begin{sidewaysfigure}
\centering
\begin{tikzpicture}[
  box/.style={draw, rounded corners=4pt, minimum width=2.8cm,
    minimum height=0.8cm, align=center, font=\small},
  source/.style={box, fill=blue!15},
  iqt/.style={box, fill=orange!25, minimum width=3.5cm, minimum height=1cm,
    font=\small\bfseries},
  target/.style={box, fill=green!20},
  arr/.style={-Stealth, thick},
  dbl/.style={Stealth-Stealth, thick, dashed},
]
  % Quellsprachen
  \node[source] (py)   at (0,4)   {Python\\SQLAlchemy};
  \node[source] (java) at (0,2)   {Java\\Streams / JPA};
  \node[source] (cs)   at (0,0)   {C\# / LINQ};

  % Parser
  \node[box, fill=yellow!20] (pp) at (3.5,4) {Python\\Parser};
  \node[box, fill=yellow!20] (jp) at (3.5,2) {Java\\Parser};
  \node[box, fill=yellow!20] (cp) at (3.5,0) {C\#\\Parser};

  % IQT
  \node[iqt] (iqt) at (7,2) {IQT\\Intermediate\\Query Tree};

  % Normalisierung
  \node[box, fill=purple!15] (norm) at (10.5,2) {IQT\\Normalisierer};

  % SQL-Generatoren
  \node[target] (ansi)  at (14,4)   {ANSI SQL:2016};
  \node[target] (pg)    at (14,2.7) {PostgreSQL};
  \node[target] (mysql) at (14,1.4) {MySQL 8};
  \node[target] (ms)    at (14,0.1) {MSSQL / Oracle};

  % Pfeile
  \draw[arr] (py)   -- (pp);
  \draw[arr] (java) -- (jp);
  \draw[arr] (cs)   -- (cp);
  \draw[arr] (pp)   -- (iqt);
  \draw[arr] (jp)   -- (iqt);
  \draw[arr] (cp)   -- (iqt);
  \draw[arr] (iqt)  -- (norm);
  \draw[arr] (norm) -- (ansi);
  \draw[arr] (norm) -- (pg);
  \draw[arr] (norm) -- (mysql);
  \draw[arr] (norm) -- (ms);

  % Beschriftung
  \node[font=\tiny, text=gray] at (1.7,4.3)  {AST-Extraktion};
  \node[font=\tiny, text=gray] at (1.7,2.3)  {AST-Extraktion};
  \node[font=\tiny, text=gray] at (1.7,0.3)  {AST-Extraktion};
  \node[font=\tiny, text=gray] at (5.2,3)    {IQT-Aufbau};
  \node[font=\tiny, text=gray] at (8.7,2.3)  {Normalisierung};
  \node[font=\tiny, text=gray] at (12.2,3)   {Codegen};
\end{tikzpicture}
\caption{UQTL-Übersetzungspipeline: von Quellsprachen über IQT zum SQL-Zieldialekt}
\label{fig:pipeline}
\end{sidewaysfigure}

\subsection{IQT-Knotentypen}
\label{subsec:iqt_knoten}

Der IQT ist ein gerichteter azyklischer Baum (DAG) aus typisierten Knoten.
Jeder Knoten entspricht einem relationalen oder logischen Operator.

\begin{definition}[IQT-Knotentypen]
\label{def:iqt_knoten}
Die Menge der zulässigen IQT-Knotentypen $\mathcal{N}$ ist definiert als:
\begin{align*}
\mathcal{N} = \{
  &\mathtt{SCAN},\; \mathtt{FILTER},\; \mathtt{PROJECT},\; \\
  &\mathtt{JOIN},\; \mathtt{GROUP\_AGG},\; \mathtt{SORT},\; \\
  &\mathtt{LIMIT},\; \mathtt{WINDOW},\; \mathtt{SUBQUERY},\; \\
  &\mathtt{SET\_OP},\; \mathtt{WITH}\;\}
\end{align*}
\end{definition}

\subsubsection{Knoten-Spezifikation im Detail}

\textbf{SCAN-Knoten}
Der \texttt{SCAN}-Knoten repräsentiert den Basistabellenzugriff:

\begin{lstlisting}[style=sqlstyle, caption={IQT-Pseudonotation: SCAN},
  label=lst:iqt_scan]
SCAN {
  table    : QualifiedName   -- Schema + Tabellenname
  alias    : Identifier?     -- optionaler Alias
  hint     : IndexHint?      -- optionaler Index-Hint (Erweiterung)
}
\end{lstlisting}

\textbf{FILTER-Knoten}
\begin{lstlisting}[style=sqlstyle, caption={IQT-Pseudonotation: FILTER},
  label=lst:iqt_filter]
FILTER {
  source   : IQTNode         -- Eingaberelation
  predicate: Expr            -- boolescher Ausdruck
  position : {WHERE | HAVING}
}

Expr ::=
  | CompareExpr { left: Expr, op: {EQ|NEQ|LT|GT|LTE|GTE}, right: Expr }
  | LogicExpr   { op: {AND|OR|NOT}, operands: Expr+ }
  | InExpr      { value: Expr, list: Expr+ | SubQuery }
  | BetweenExpr { value: Expr, lo: Expr, hi: Expr }
  | LikeExpr    { value: Expr, pattern: StringLiteral, escape: Char? }
  | NullExpr    { value: Expr, negated: Bool }
  | ExistsExpr  { subquery: SubQuery, negated: Bool }
  | Literal     { type: DataType, value: Any }
  | ColumnRef   { table: Identifier?, column: Identifier }
  | FuncCall    { name: QualifiedName, args: Expr*, distinct: Bool }
\end{lstlisting}

\textbf{PROJECT-Knoten}
\begin{lstlisting}[style=sqlstyle, caption={IQT-Pseudonotation: PROJECT},
  label=lst:iqt_project]
PROJECT {
  source    : IQTNode
  columns   : ProjectionItem+
  distinct  : Bool
}

ProjectionItem ::=
  | Star       {}
  | Expression { expr: Expr, alias: Identifier? }
\end{lstlisting}

\textbf{JOIN-Knoten}
\begin{lstlisting}[style=sqlstyle, caption={IQT-Pseudonotation: JOIN},
  label=lst:iqt_join]
JOIN {
  left      : IQTNode
  right     : IQTNode
  type      : {INNER | LEFT_OUTER | RIGHT_OUTER | FULL_OUTER | CROSS}
  condition : JoinCondition
}

JoinCondition ::=
  | OnCondition    { predicate: Expr }
  | UsingCondition { columns: Identifier+ }
  | Natural        {}
\end{lstlisting}

\textbf{GROUP\_AGG-Knoten}
\begin{lstlisting}[style=sqlstyle, caption={IQT-Pseudonotation: GROUP\_AGG},
  label=lst:iqt_group]
GROUP_AGG {
  source     : IQTNode
  groupBy    : Expr*           -- leer = Gesamt-Aggregation
  aggregates : AggItem+
  having     : Expr?
}

AggItem ::= {
  func  : {COUNT|SUM|AVG|MIN|MAX|STDDEV|VAR|STRING_AGG|...}
  arg   : Expr          -- COUNT(*) => StarExpr
  distinct: Bool
  alias : Identifier?
}
\end{lstlisting}

\textbf{SORT-Knoten}
\begin{lstlisting}[style=sqlstyle, caption={IQT-Pseudonotation: SORT},
  label=lst:iqt_sort]
SORT {
  source  : IQTNode
  items   : SortItem+
}

SortItem ::= {
  expr      : Expr
  direction : {ASC | DESC}
  nulls     : {NULLS_FIRST | NULLS_LAST | DEFAULT}
}
\end{lstlisting}

\textbf{LIMIT-Knoten}
\begin{lstlisting}[style=sqlstyle, caption={IQT-Pseudonotation: LIMIT},
  label=lst:iqt_limit]
LIMIT {
  source : IQTNode
  count  : Expr | Null    -- NULL = kein Limit
  offset : Expr | Null    -- NULL = kein Offset
}
\end{lstlisting}

\textbf{WINDOW-Knoten}
\begin{lstlisting}[style=sqlstyle, caption={IQT-Pseudonotation: WINDOW},
  label=lst:iqt_window]
WINDOW {
  source       : IQTNode
  definitions  : WindowDef+
}

WindowDef ::= {
  alias        : Identifier
  partitionBy  : Expr*
  orderBy      : SortItem*
  frame        : WindowFrame?
  func         : WindowFunc
  resultAlias  : Identifier
}

WindowFrame ::= {
  unit  : {ROWS | RANGE}
  start : FrameBound
  end   : FrameBound
}

FrameBound ::=
  | UNBOUNDED_PRECEDING | CURRENT_ROW | UNBOUNDED_FOLLOWING
  | Preceding { n: Expr } | Following { n: Expr }
\end{lstlisting}

\subsection{Normalisierungsregeln des IQT}
\label{subsec:iqt_norm}

Um die eindeutige SQL-Ausgabe zu garantieren, werden IQT-Bäume vor der
SQL-Generierung normalisiert. Die Normalisierung umfasst folgende Schritte:

\begin{regel}[Prädikat-Normalisierung]
\label{regel:pred_norm}
Zusammengesetzte \texttt{AND}-Prädikate werden in einer kanonischen
Reihenfolge sortiert (lexikografisch nach Spaltenname, dann Operator):
\[
  p_1 \wedge p_2 \wedge \ldots \wedge p_n \;\rightarrow\;
  p_{\sigma(1)} \wedge p_{\sigma(2)} \wedge \ldots \wedge p_{\sigma(n)}
\]
mit $\sigma$ als lexikografischer Permutation.
\end{regel}

\begin{regel}[Alias-Normalisierung]
\label{regel:alias_norm}
Generierte Tabellenaliase folgen dem Schema \texttt{t1, t2, \ldots} in
der Reihenfolge des Auftretens im Baum (Pre-Order-Traversal).
\end{regel}

\begin{regel}[Boolean-Vereinfachung]
\label{regel:bool_simp}
Redundante boolesche Ausdrücke werden eliminiert:
\begin{align*}
  p \wedge \mathtt{TRUE} &\;\rightarrow\; p \\
  p \vee \mathtt{FALSE}  &\;\rightarrow\; p \\
  \neg\neg p             &\;\rightarrow\; p \\
  p \wedge p             &\;\rightarrow\; p
\end{align*}
\end{regel}

\begin{regel}[SELECT-Reihenfolge]
\label{regel:select_order}
Explizit genannte Projekti\-ons\-spalten werden in der Reihenfolge ihrer
Deklaration ausgegeben. \texttt{SELECT *} wird nur bei echten Stern-Projektionen
verwendet; implizite Stern-Projektionen werden beibehalten.
\end{regel}

\begin{regel}[JOIN-Kanonisierung]
\label{regel:join_canon}
Mehrfach-Joins werden left-assoziativ normiert.
\texttt{CROSS JOIN} ohne \texttt{ON}-Klausel wird nur verwendet, wenn
explizit kein Prädikat vorliegt; sonst wird \texttt{INNER JOIN ... ON}
bevorzugt.
\end{regel}

\newpage
% ═══════════════════════════════════════════════════════════════════════════════
\section{Übersetzungsregeln und Operatoren}
\label{sec:operatoren}
% ═══════════════════════════════════════════════════════════════════════════════

Dieses Kapitel definiert die vollständigen Übersetzungsregeln von den
Quellsprachen Python, Java und C\# in den IQT sowie vom IQT nach ANSI SQL.
Die Regeln werden als Transformationstabellen und formale Schreibweisen
dargestellt.

% ─────────────────────────────────────────────────────────────────────────────
\subsection{Operator 1: Filterung (WHERE / HAVING)}
\label{subsec:op_filter}
% ─────────────────────────────────────────────────────────────────────────────

\begin{sidewaystable}
\centering
\caption{Übersetzungstabelle: Filteroperatoren}
\label{tab:filter_ops}
\renewcommand{\arraystretch}{1.4}
\begin{tabularx}{\textwidth}{lXXX}
\toprule
\textbf{IQT-Prädikat} & \textbf{Python} & \textbf{Java} & \textbf{C\#} \\
\midrule
$\mathtt{EQ}(a, v)$
  & \texttt{Model.col == v}
  & \texttt{cb.equal(r.get("col"),v)}
  & \texttt{x.Col == v} \\
$\mathtt{NEQ}(a, v)$
  & \texttt{Model.col != v}
  & \texttt{cb.notEqual(...)}
  & \texttt{x.Col != v} \\
$\mathtt{GT}(a, v)$
  & \texttt{Model.col > v}
  & \texttt{cb.gt(...)}
  & \texttt{x.Col > v} \\
$\mathtt{GTE}(a, v)$
  & \texttt{Model.col >= v}
  & \texttt{cb.ge(...)}
  & \texttt{x.Col >= v} \\
$\mathtt{IN}(a, list)$
  & \texttt{Model.col.in\_(list)}
  & \texttt{root.get("col").in(list)}
  & \texttt{list.Contains(x.Col)} \\
$\mathtt{BETWEEN}$
  & \texttt{col.between(a,b)}
  & \texttt{cb.between(r.get("col"),a,b)}
  & \texttt{x.Col >= a \&\& x.Col <= b} \\
$\mathtt{LIKE}(a, p)$
  & \texttt{Model.col.like(p)}
  & \texttt{cb.like(r.get("col"),p)}
  & \texttt{EF.Functions.Like(x.Col,p)} \\
$\mathtt{IS\_NULL}$
  & \texttt{Model.col == None}
  & \texttt{cb.isNull(r.get("col"))}
  & \texttt{x.Col == null} \\
$\mathtt{AND}$
  & \texttt{\&} oder \texttt{and\_()}
  & \texttt{cb.and(...)}
  & \texttt{\&\&} \\
$\mathtt{OR}$
  & \texttt{|} oder \texttt{or\_()}
  & \texttt{cb.or(...)}
  & \texttt{||} \\
\bottomrule
\end{tabularx}
\end{sidewaystable}

\textbf{SQL-Generierungsregel für FILTER:}
\begin{align*}
\mathrm{SQL}(\mathtt{FILTER}(S, p, \mathtt{WHERE}))
  &= \mathrm{SQL}(S) + \texttt{ WHERE } + \mathrm{SQL}(p) \\
\mathrm{SQL}(\mathtt{FILTER}(S, p, \mathtt{HAVING}))
  &= \mathrm{SQL}(S) + \texttt{ HAVING } + \mathrm{SQL}(p)
\end{align*}

% ─────────────────────────────────────────────────────────────────────────────
\subsection{Operator 2: Projektion (SELECT)}
\label{subsec:op_project}
% ─────────────────────────────────────────────────────────────────────────────

\begin{sidewaystable}
\centering
\caption{Übersetzungstabelle: Projektionsoperatoren}
\label{tab:project_ops}
\renewcommand{\arraystretch}{1.4}
\begin{tabularx}{\textwidth}{lXXX}
\toprule
\textbf{Operation} & \textbf{Python} & \textbf{Java} & \textbf{C\#} \\
\midrule
Alle Spalten
  & \texttt{select(Model)}
  & \texttt{cq.select(root)}
  & \texttt{.Select(x => x)} \\
Einzelne Spalten
  & \texttt{select(M.a, M.b)}
  & \texttt{cb.construct(DTO, r.get("a"), r.get("b"))}
  & \texttt{.Select(x => new\{x.A,x.B\})} \\
Ausdruck
  & \texttt{select(func.upper(M.name))}
  & \texttt{cb.upper(root.get("name"))}
  & \texttt{.Select(x => x.Name.ToUpper())} \\
Distinct
  & \texttt{select(M.col).distinct()}
  & \texttt{cq.distinct(true)}
  & \texttt{.Select(x => x.Col).Distinct()} \\
\bottomrule
\end{tabularx}
\end{sidewaystable}

% ─────────────────────────────────────────────────────────────────────────────
\subsection{Operator 3: Sortierung (ORDER BY)}
\label{subsec:op_sort}
% ─────────────────────────────────────────────────────────────────────────────

\begin{sidewaystable}
\centering
\caption{Übersetzungstabelle: Sortieroperatoren}
\label{tab:sort_ops}
\renewcommand{\arraystretch}{1.4}
\begin{tabularx}{\textwidth}{lXXX}
\toprule
\textbf{IQT-Operation} & \textbf{Python} & \textbf{Java} & \textbf{C\#} \\
\midrule
ASC
  & \texttt{order\_by(M.col)}
  & \texttt{cb.asc(r.get("col"))}
  & \texttt{OrderBy(x => x.Col)} \\
DESC
  & \texttt{order\_by(desc(M.col))}
  & \texttt{cb.desc(r.get("col"))}
  & \texttt{OrderByDescending(x => x.Col)} \\
Multi-Sort
  & \texttt{order\_by(M.a, desc(M.b))}
  & \texttt{.orderBy(cb.asc(...), cb.desc(...))}
  & \texttt{.OrderBy(x => x.A).ThenByDescending(x => x.B)} \\
NULLS FIRST
  & \texttt{nullsfirst(M.col)}
  & \texttt{cb.asc(r.get("col")).nullsFirst()}
  & \texttt{EF nicht nativ} \\
\bottomrule
\end{tabularx}
\end{sidewaystable}

\textbf{SQL-Generierungsregel:}
\[
\mathrm{SQL}(\mathtt{SORT}(S, [s_1, \ldots, s_n])) =
  \mathrm{SQL}(S) + \texttt{ ORDER BY } +
  \mathrm{SQL}(s_1) + \texttt{, } + \ldots + \texttt{, } + \mathrm{SQL}(s_n)
\]
\[
\mathrm{SQL}(\mathtt{SortItem}(e, \mathtt{ASC}, \mathtt{DEFAULT})) =
  \mathrm{SQL}(e) + \texttt{ ASC}
\]

% ─────────────────────────────────────────────────────────────────────────────
\subsection{Operator 4: Begrenzung (LIMIT / OFFSET)}
\label{subsec:op_limit}
% ─────────────────────────────────────────────────────────────────────────────

\begin{sidewaystable}
\centering
\caption{Übersetzungstabelle: LIMIT / OFFSET je Datenbankdialekt}
\label{tab:limit_ops}
\renewcommand{\arraystretch}{1.4}
\begin{tabularx}{\textwidth}{lXX}
\toprule
\textbf{Sprache / Dialekt} & \textbf{Quellausdruck} & \textbf{SQL-Ausgabe} \\
\midrule
Python (SA)
  & \texttt{.limit(10).offset(20)}
  & \texttt{LIMIT 10 OFFSET 20} \\
Java (JPA)
  & \texttt{.setMaxResults(10).setFirstResult(20)}
  & \texttt{LIMIT 10 OFFSET 20} \\
C\# (LINQ)
  & \texttt{.Take(10).Skip(20)}
  & \texttt{LIMIT 10 OFFSET 20} \\
\midrule
\multicolumn{3}{l}{\textit{Dialektspezifische Ausgabe (Erweiterungsmodus):}} \\
MSSQL
  & -- (wie oben)
  & \texttt{OFFSET 20 ROWS FETCH NEXT 10 ROWS ONLY} \\
Oracle $<$12c
  & -- (wie oben)
  & \texttt{WHERE ROWNUM BETWEEN 21 AND 30} (umgeschrieben) \\
\bottomrule
\end{tabularx}
\end{sidewaystable}

% ─────────────────────────────────────────────────────────────────────────────
\subsection{Operator 5: Aggregation (GROUP BY)}
\label{subsec:op_agg}
% ─────────────────────────────────────────────────────────────────────────────

\begin{sidewaystable}
\centering
\caption{Übersetzungstabelle: Aggregationsfunktionen}
\label{tab:agg_ops}
\renewcommand{\arraystretch}{1.4}
\begin{tabularx}{\textwidth}{lXXX}
\toprule
\textbf{IQT-Funktion} & \textbf{Python} & \textbf{Java} & \textbf{C\#} \\
\midrule
\texttt{COUNT(*)}
  & \texttt{func.count('*')}
  & \texttt{cb.count(root)}
  & \texttt{.Count()} \\
\texttt{COUNT(col)}
  & \texttt{func.count(M.col)}
  & \texttt{cb.count(r.get("col"))}
  & \texttt{.Count(x => x.Col != null)} \\
\texttt{SUM(col)}
  & \texttt{func.sum(M.col)}
  & \texttt{cb.sum(r.get("col"))}
  & \texttt{.Sum(x => x.Col)} \\
\texttt{AVG(col)}
  & \texttt{func.avg(M.col)}
  & \texttt{cb.avg(r.get("col"))}
  & \texttt{.Average(x => x.Col)} \\
\texttt{MIN(col)}
  & \texttt{func.min(M.col)}
  & \texttt{cb.min(r.get("col"))}
  & \texttt{.Min(x => x.Col)} \\
\texttt{MAX(col)}
  & \texttt{func.max(M.col)}
  & \texttt{cb.max(r.get("col"))}
  & \texttt{.Max(x => x.Col)} \\
GROUP BY
  & \texttt{group\_by(M.dep)}
  & \texttt{cq.groupBy(r.get("dep"))}
  & \texttt{.GroupBy(x => x.Dep)} \\
HAVING
  & \texttt{having(func.count('*') > 5)}
  & \texttt{cq.having(cb.gt(cb.count(r),5))}
  & \texttt{.Where(g => g.Count() > 5)} \\
\bottomrule
\end{tabularx}
\end{sidewaystable}

% ─────────────────────────────────────────────────────────────────────────────
\subsection{Operator 6: Joins}
\label{subsec:op_join}
% ─────────────────────────────────────────────────────────────────────────────

\begin{sidewaystable}
\centering
\caption{Übersetzungstabelle: Join-Typen}
\label{tab:join_ops}
\renewcommand{\arraystretch}{1.4}
\begin{tabularx}{\textwidth}{lXXX}
\toprule
\textbf{Join-Typ} & \textbf{Python} & \textbf{Java} & \textbf{C\#} \\
\midrule
INNER JOIN
  & \texttt{join(Order, User.id == Order.user\_id)}
  & \texttt{root.join("orders")}
  & \texttt{join in context.Orders on u.Id equals o.UserId} \\
LEFT OUTER
  & \texttt{outerjoin(Order, ...)}
  & \texttt{root.join(..., JoinType.LEFT)}
  & \texttt{join o in ... into grp from o in grp.DefaultIfEmpty()} \\
RIGHT OUTER
  & \texttt{(SA kein nativer Support)}
  & \texttt{JoinType.RIGHT}
  & \texttt{(symmetrisch zu LEFT)} \\
CROSS JOIN
  & \texttt{select\_from(A, B)}
  & \texttt{cq.from(A.class); cq.from(B.class)}
  & \texttt{from a in As from b in Bs select ...} \\
\bottomrule
\end{tabularx}
\end{sidewaystable}

\begin{regel}[JOIN-Übersetzung]
\label{regel:join}
Für einen \texttt{INNER JOIN} zwischen Relation $L$ (Alias $t_1$) und $R$ (Alias $t_2$)
mit Prädikat $p$ ergibt sich:
\[
\mathrm{SQL}(\mathtt{JOIN}(L, R, \mathtt{INNER}, p)) =
  \mathrm{SQL}(L) + \texttt{ INNER JOIN } + \mathrm{SQL}(R) +
  \texttt{ ON } + \mathrm{SQL}(p)
\]
\end{regel}

% ─────────────────────────────────────────────────────────────────────────────
\subsection{Operator 7: Fensterfunktionen (OVER)}
\label{subsec:op_window}
% ─────────────────────────────────────────────────────────────────────────────

Fensterfunktionen (Window Functions, ANSI SQL:2003+) werden in UQTL ab
Konformitätsstufe CL-3 unterstützt.

\begin{sidewaystable}
\centering
\caption{Übersetzungstabelle: Fensterfunktionen}
\label{tab:window_ops}
\renewcommand{\arraystretch}{1.4}
\begin{tabularx}{\textwidth}{lXXX}
\toprule
\textbf{SQL-Funktion} & \textbf{Python} & \textbf{Java} & \textbf{C\#} \\
\midrule
\texttt{ROW\_NUMBER()}
  & \texttt{func.row\_number().over(...)}
  & \texttt{cb.rowNumber().over(...)}
  & \texttt{EF.Functions.RowNumber(...)} \\
\texttt{RANK()}
  & \texttt{func.rank().over(...)}
  & \texttt{cb.rank().over(...)}
  & \texttt{EF.Functions.Rank(...)} \\
\texttt{DENSE\_RANK()}
  & \texttt{func.dense\_rank().over(...)}
  & \texttt{cb.denseRank().over(...)}
  & \texttt{EF.Functions.DenseRank(...)} \\
\texttt{LAG(col, n)}
  & \texttt{func.lag(M.col, n).over(...)}
  & \texttt{cb.lag(r.get("col"), n)}
  & \texttt{EF.Functions.Lag(x.Col, n)} \\
\texttt{LEAD(col, n)}
  & \texttt{func.lead(M.col, n).over(...)}
  & \texttt{cb.lead(r.get("col"), n)}
  & \texttt{EF.Functions.Lead(x.Col, n)} \\
\texttt{SUM() OVER}
  & \texttt{func.sum(M.col).over(...)}
  & \texttt{cb.sum(r.get("col")).over(...)}
  & \texttt{.Sum(x.Col).Over(...)} \\
\bottomrule
\end{tabularx}
\end{sidewaystable}

% ─────────────────────────────────────────────────────────────────────────────
\subsection{Operator 8: Unterabfragen und CTEs}
\label{subsec:op_subquery}
% ─────────────────────────────────────────────────────────────────────────────

\begin{sidewaystable}
\centering
\caption{Übersetzungstabelle: Unterabfragen}
\label{tab:subquery_ops}
\renewcommand{\arraystretch}{1.4}
\begin{tabularx}{\textwidth}{lXXX}
\toprule
\textbf{Typ} & \textbf{Python} & \textbf{Java} & \textbf{C\#} \\
\midrule
Scalar Subquery
  & \texttt{subquery()}
  & \texttt{cb.subquery(Long.class)}
  & \texttt{context.Baz.Where(...).Select(x => x.Val).FirstOrDefault()} \\
EXISTS
  & \texttt{exists(subq)}
  & \texttt{cb.exists(subq)}
  & \texttt{.Any(cond)} \\
IN (Subquery)
  & \texttt{M.col.in\_(subq)}
  & \texttt{root.get("col").in(subq)}
  & \texttt{list.Contains(x.Col)} (mit AsQueryable) \\
CTE (WITH)
  & \texttt{cte = subq.cte("name")}
  & \texttt{(JPA kein nativer CTE)}
  & \texttt{(EF Core $\geq$7: FromSqlRaw)} \\
\bottomrule
\end{tabularx}
\end{sidewaystable}

% ─────────────────────────────────────────────────────────────────────────────
\subsection{Operator 9: Mengenoperationen}
\label{subsec:op_set}
% ─────────────────────────────────────────────────────────────────────────────

\begin{table}[H]
\centering
\caption{Übersetzungstabelle: Mengenoperationen}
\label{tab:set_ops}
\renewcommand{\arraystretch}{1.4}
\begin{tabularx}{\textwidth}{lXXX}
\toprule
\textbf{IQT} & \textbf{Python} & \textbf{Java} & \textbf{C\#} \\
\midrule
UNION
  & \texttt{q1.union(q2)}
  & \texttt{(via CriteriaQuery)} 
  & \texttt{q1.Union(q2)} \\
UNION ALL
  & \texttt{q1.union\_all(q2)}
  & \texttt{(JPA: JPQL UNION)} 
  & \texttt{q1.Concat(q2)} \\
INTERSECT
  & \texttt{q1.intersect(q2)}
  & \texttt{(selten nativ)}
  & \texttt{q1.Intersect(q2)} \\
EXCEPT
  & \texttt{q1.except\_(q2)}
  & \texttt{(selten nativ)}
  & \texttt{q1.Except(q2)} \\
\bottomrule
\end{tabularx}
\end{table}

\newpage
% ═══════════════════════════════════════════════════════════════════════════════
\section{Praxisbeispiele}
\label{sec:beispiele}
% ═══════════════════════════════════════════════════════════════════════════════

Dieses Kapitel demonstriert den UQTL-Standard anhand vollständiger,
realistischer Datenbankszenarien. Für jedes Beispiel wird zunächst das
Datenbankschema beschrieben, dann die Abfrage in allen drei Quellsprachen
formuliert, und schließlich das vom Standard definierte normalisierte
SQL-Ergebnis angegeben.

\subsection{Datenmodell für die Beispiele}
\label{subsec:datenmodell}

Alle Beispiele basieren auf folgendem vereinfachten E-Commerce-Schema:

\begin{lstlisting}[style=sqlstyle, caption={Datenbankschema für Beispiele},
  label=lst:schema]
CREATE TABLE customers (
    id          INTEGER PRIMARY KEY,
    name        VARCHAR(200) NOT NULL,
    email       VARCHAR(200) UNIQUE,
    country     VARCHAR(50),
    active      BOOLEAN NOT NULL DEFAULT TRUE,
    created_at  TIMESTAMP NOT NULL,
    tier        VARCHAR(20)  -- 'bronze', 'silver', 'gold', 'platinum'
);

CREATE TABLE orders (
    id          INTEGER PRIMARY KEY,
    customer_id INTEGER REFERENCES customers(id),
    order_date  DATE NOT NULL,
    status      VARCHAR(30),  -- 'pending','confirmed','shipped','delivered'
    total       DECIMAL(12,2) NOT NULL
);

CREATE TABLE order_items (
    id          INTEGER PRIMARY KEY,
    order_id    INTEGER REFERENCES orders(id),
    product_id  INTEGER REFERENCES products(id),
    quantity    INTEGER NOT NULL,
    unit_price  DECIMAL(10,2) NOT NULL
);

CREATE TABLE products (
    id          INTEGER PRIMARY KEY,
    name        VARCHAR(200) NOT NULL,
    category    VARCHAR(100),
    price       DECIMAL(10,2) NOT NULL,
    stock       INTEGER NOT NULL DEFAULT 0
);
\end{lstlisting}

% ─────────────────────────────────────────────────────────────────────────────
\subsection{Beispiel 1: Einfache Filterung und Projektion}
\label{subsec:bsp1}
% ─────────────────────────────────────────────────────────────────────────────

\textbf{Aufgabe:} Alle aktiven Kunden aus Deutschland, nur Name und E-Mail,
alphabetisch sortiert.

\begin{lstlisting}[style=pythonstyle, caption={Beispiel 1 -- Python (SQLAlchemy 2.0)},
  label=lst:bsp1_py]
from sqlalchemy import select
from models import Customer

stmt = (
    select(Customer.name, Customer.email)
    .where(Customer.active == True)
    .where(Customer.country == "DE")
    .order_by(Customer.name.asc())
)
results = session.execute(stmt).all()
\end{lstlisting}

\begin{lstlisting}[style=javastyle, caption={Beispiel 1 -- Java (JPA Criteria API)},
  label=lst:bsp1_java]
CriteriaBuilder cb = em.getCriteriaBuilder();
CriteriaQuery<Object[]> cq = cb.createQuery(Object[].class);
Root<Customer> c = cq.from(Customer.class);

cq.multiselect(c.get("name"), c.get("email"))
  .where(
      cb.and(
          cb.isTrue(c.get("active")),
          cb.equal(c.get("country"), "DE")
      )
  )
  .orderBy(cb.asc(c.get("name")));

List<Object[]> results = em.createQuery(cq).getResultList();
\end{lstlisting}

\begin{lstlisting}[style=csharpstyle, caption={Beispiel 1 -- C\# (EF Core / LINQ)},
  label=lst:bsp1_cs]
// Query Syntax
var results = from c in context.Customers
              where c.Active == true && c.Country == "DE"
              orderby c.Name ascending
              select new { c.Name, c.Email };

// Equivalent Method Syntax
var results = context.Customers
    .Where(c => c.Active == true && c.Country == "DE")
    .OrderBy(c => c.Name)
    .Select(c => new { c.Name, c.Email });
\end{lstlisting}

\begin{successbox}
\textbf{UQTL-Normiertes SQL (identisch für alle drei Quellsprachen):}
\begin{lstlisting}[style=sqlstyle, numbers=none]
SELECT t1.name, t1.email
FROM customers AS t1
WHERE t1.active = TRUE
  AND t1.country = 'DE'
ORDER BY t1.name ASC
\end{lstlisting}
\end{successbox}

% ─────────────────────────────────────────────────────────────────────────────
\subsection{Beispiel 2: Aggregation mit GROUP BY und HAVING}
\label{subsec:bsp2}
% ─────────────────────────────────────────────────────────────────────────────

\textbf{Aufgabe:} Anzahl bestätigter Bestellungen und Gesamtumsatz pro Kunde,
nur Kunden mit mehr als 3 Bestellungen, absteigend nach Umsatz.

\begin{lstlisting}[style=pythonstyle, caption={Beispiel 2 -- Python},
  label=lst:bsp2_py]
from sqlalchemy import select, func

stmt = (
    select(
        Order.customer_id,
        func.count(Order.id).label("order_count"),
        func.sum(Order.total).label("total_revenue")
    )
    .where(Order.status == "confirmed")
    .group_by(Order.customer_id)
    .having(func.count(Order.id) > 3)
    .order_by(func.sum(Order.total).desc())
)
\end{lstlisting}

\begin{lstlisting}[style=javastyle, caption={Beispiel 2 -- Java},
  label=lst:bsp2_java]
CriteriaBuilder cb = em.getCriteriaBuilder();
CriteriaQuery<Object[]> cq = cb.createQuery(Object[].class);
Root<Order> o = cq.from(Order.class);

Expression<Long>    cnt = cb.count(o.get("id"));
Expression<Double>  rev = cb.sum(o.<Double>get("total"));

cq.multiselect(o.get("customerId"), cnt, rev)
  .where(cb.equal(o.get("status"), "confirmed"))
  .groupBy(o.get("customerId"))
  .having(cb.gt(cnt, 3L))
  .orderBy(cb.desc(rev));
\end{lstlisting}

\begin{lstlisting}[style=csharpstyle, caption={Beispiel 2 -- C\#},
  label=lst:bsp2_cs]
var results = context.Orders
    .Where(o => o.Status == "confirmed")
    .GroupBy(o => o.CustomerId)
    .Where(g => g.Count() > 3)
    .OrderByDescending(g => g.Sum(o => o.Total))
    .Select(g => new {
        CustomerId   = g.Key,
        OrderCount   = g.Count(),
        TotalRevenue = g.Sum(o => o.Total)
    });
\end{lstlisting}

\begin{successbox}
\textbf{UQTL-Normiertes SQL:}
\begin{lstlisting}[style=sqlstyle, numbers=none]
SELECT t1.customer_id,
       COUNT(t1.id)    AS order_count,
       SUM(t1.total)   AS total_revenue
FROM orders AS t1
WHERE t1.status = 'confirmed'
GROUP BY t1.customer_id
HAVING COUNT(t1.id) > 3
ORDER BY SUM(t1.total) DESC
\end{lstlisting}
\end{successbox}

% ─────────────────────────────────────────────────────────────────────────────
\subsection{Beispiel 3: INNER JOIN mit mehreren Bedingungen}
\label{subsec:bsp3}
% ─────────────────────────────────────────────────────────────────────────────

\textbf{Aufgabe:} Name und E-Mail der Kunden, die mindestens eine
gelieferte Bestellung im Wert über 500~EUR hatten, im Jahr 2024.

\begin{lstlisting}[style=pythonstyle, caption={Beispiel 3 -- Python (JOIN)},
  label=lst:bsp3_py]
from sqlalchemy import select, extract

stmt = (
    select(Customer.name, Customer.email)
    .join(Order, Customer.id == Order.customer_id)
    .where(Order.status == "delivered")
    .where(Order.total > 500)
    .where(extract("year", Order.order_date) == 2024)
    .distinct()
)
\end{lstlisting}

\begin{lstlisting}[style=javastyle, caption={Beispiel 3 -- Java (JOIN)},
  label=lst:bsp3_java]
CriteriaBuilder cb = em.getCriteriaBuilder();
CriteriaQuery<Object[]> cq = cb.createQuery(Object[].class);
Root<Customer> c = cq.from(Customer.class);
Join<Customer, Order> o = c.join("orders", JoinType.INNER);

cq.multiselect(c.get("name"), c.get("email"))
  .where(cb.and(
      cb.equal(o.get("status"), "delivered"),
      cb.gt(o.<BigDecimal>get("total"), new BigDecimal("500")),
      cb.equal(
          cb.function("YEAR", Integer.class, o.get("orderDate")),
          2024
      )
  ))
  .distinct(true);
\end{lstlisting}

\begin{lstlisting}[style=csharpstyle, caption={Beispiel 3 -- C\# (JOIN)},
  label=lst:bsp3_cs]
var results = (from c in context.Customers
               join o in context.Orders
                 on c.Id equals o.CustomerId
               where o.Status == "delivered"
                  && o.Total > 500m
                  && o.OrderDate.Year == 2024
               select new { c.Name, c.Email })
              .Distinct();
\end{lstlisting}

\begin{successbox}
\textbf{UQTL-Normiertes SQL:}
\begin{lstlisting}[style=sqlstyle, numbers=none]
SELECT DISTINCT t1.name, t1.email
FROM customers AS t1
INNER JOIN orders AS t2
  ON t1.id = t2.customer_id
WHERE t2.status = 'delivered'
  AND t2.total > 500
  AND EXTRACT(YEAR FROM t2.order_date) = 2024
\end{lstlisting}
\end{successbox}

% ─────────────────────────────────────────────────────────────────────────────
\subsection{Beispiel 4: LEFT OUTER JOIN und NULL-Behandlung}
\label{subsec:bsp4}
% ─────────────────────────────────────────────────────────────────────────────

\textbf{Aufgabe:} Alle Kunden mit Anzahl ihrer Bestellungen (auch solche ohne Bestellungen).

\begin{lstlisting}[style=pythonstyle, caption={Beispiel 4 -- Python (LEFT JOIN)},
  label=lst:bsp4_py]
from sqlalchemy import select, func, outerjoin

stmt = (
    select(
        Customer.id,
        Customer.name,
        func.count(Order.id).label("order_count")
    )
    .outerjoin(Order, Customer.id == Order.customer_id)
    .group_by(Customer.id, Customer.name)
    .order_by(Customer.name)
)
\end{lstlisting}

\begin{lstlisting}[style=javastyle, caption={Beispiel 4 -- Java (LEFT JOIN)},
  label=lst:bsp4_java]
CriteriaBuilder cb = em.getCriteriaBuilder();
CriteriaQuery<Object[]> cq = cb.createQuery(Object[].class);
Root<Customer> c = cq.from(Customer.class);
Join<Customer, Order> o = c.join("orders", JoinType.LEFT);

cq.multiselect(c.get("id"), c.get("name"), cb.count(o.get("id")))
  .groupBy(c.get("id"), c.get("name"))
  .orderBy(cb.asc(c.get("name")));
\end{lstlisting}

\begin{lstlisting}[style=csharpstyle, caption={Beispiel 4 -- C\# (LEFT JOIN)},
  label=lst:bsp4_cs]
var results = from c in context.Customers
              join o in context.Orders
                on c.Id equals o.CustomerId into orders
              from o in orders.DefaultIfEmpty()
              group o by new { c.Id, c.Name } into g
              orderby g.Key.Name
              select new {
                  g.Key.Id,
                  g.Key.Name,
                  OrderCount = g.Count(o => o != null)
              };
\end{lstlisting}

\begin{successbox}
\textbf{UQTL-Normiertes SQL:}
\begin{lstlisting}[style=sqlstyle, numbers=none]
SELECT t1.id, t1.name, COUNT(t2.id) AS order_count
FROM customers AS t1
LEFT OUTER JOIN orders AS t2
  ON t1.id = t2.customer_id
GROUP BY t1.id, t1.name
ORDER BY t1.name ASC
\end{lstlisting}
\end{successbox}

% ─────────────────────────────────────────────────────────────────────────────
\subsection{Beispiel 5: Fensterfunktionen (OVER PARTITION BY)}
\label{subsec:bsp5}
% ─────────────────────────────────────────────────────────────────────────────

\textbf{Aufgabe:} Für jede Bestellung: Rang innerhalb des Kunden nach Bestellwert
(teuerste = Rang 1), laufende Summe der Bestellwerte pro Kunde.

\begin{lstlisting}[style=pythonstyle, caption={Beispiel 5 -- Python (Window Functions)},
  label=lst:bsp5_py]
from sqlalchemy import select, func
from sqlalchemy import over

rnk = func.rank().over(
    partition_by=Order.customer_id,
    order_by=Order.total.desc()
).label("rnk")

run_total = func.sum(Order.total).over(
    partition_by=Order.customer_id,
    order_by=Order.order_date.asc(),
    rows=(None, 0)  # UNBOUNDED PRECEDING to CURRENT ROW
).label("running_total")

stmt = select(
    Order.id,
    Order.customer_id,
    Order.total,
    rnk,
    run_total
).order_by(Order.customer_id, rnk)
\end{lstlisting}

\begin{lstlisting}[style=javastyle, caption={Beispiel 5 -- Java (Window via JPQL / Blaze-Persistence)},
  label=lst:bsp5_java]
// Java JPA hat keine native Window Function API.
// Blaze-Persistence oder native Queries werden benoetigt:
String jpql = """
    SELECT o.id, o.customerId, o.total,
           RANK() OVER (PARTITION BY o.customerId ORDER BY o.total DESC),
           SUM(o.total) OVER (
               PARTITION BY o.customerId
               ORDER BY o.orderDate
               ROWS BETWEEN UNBOUNDED PRECEDING AND CURRENT ROW
           )
    FROM Order o
    ORDER BY o.customerId, 4
    """;
List<Object[]> results = em.createQuery(jpql).getResultList();
\end{lstlisting}

\begin{lstlisting}[style=csharpstyle, caption={Beispiel 5 -- C\# (EF Core Window Functions)},
  label=lst:bsp5_cs]
// EF Core 8+ mit Window Function Unterstuetzung:
var results = context.Orders
    .Select(o => new {
        o.Id,
        o.CustomerId,
        o.Total,
        Rank = EF.Functions.Rank()
            .Over()
            .PartitionBy(o.CustomerId)
            .OrderByDescending(o.Total),
        RunningTotal = EF.Functions.Sum(o.Total)
            .Over()
            .PartitionBy(o.CustomerId)
            .OrderBy(o.OrderDate)
            .RowsBetween(
                EF.Functions.RowsOrRangeBoundary.UnboundedPreceding,
                EF.Functions.RowsOrRangeBoundary.CurrentRow)
    })
    .OrderBy(x => x.CustomerId)
    .ThenBy(x => x.Rank);
\end{lstlisting}

\begin{successbox}
\textbf{UQTL-Normiertes SQL:}
\begin{lstlisting}[style=sqlstyle, numbers=none]
SELECT t1.id,
       t1.customer_id,
       t1.total,
       RANK() OVER (
           PARTITION BY t1.customer_id
           ORDER BY t1.total DESC
       ) AS rnk,
       SUM(t1.total) OVER (
           PARTITION BY t1.customer_id
           ORDER BY t1.order_date ASC
           ROWS BETWEEN UNBOUNDED PRECEDING AND CURRENT ROW
       ) AS running_total
FROM orders AS t1
ORDER BY t1.customer_id ASC, rnk ASC
\end{lstlisting}
\end{successbox}

% ─────────────────────────────────────────────────────────────────────────────
\subsection{Beispiel 6: Unterabfragen (Subqueries)}
\label{subsec:bsp6}
% ─────────────────────────────────────────────────────────────────────────────

\textbf{Aufgabe:} Kunden, deren letzter Bestellwert über dem Durchschnitt
aller Bestellwerte liegt.

\begin{lstlisting}[style=pythonstyle, caption={Beispiel 6 -- Python (Subquery)},
  label=lst:bsp6_py]
from sqlalchemy import select, func

# Unterabfrage: Durchschnittlicher Bestellwert
avg_subq = select(func.avg(Order.total)).scalar_subquery()

# Unterabfrage: Letzte Bestellung je Kunde
last_order = (
    select(Order.customer_id, func.max(Order.order_date).label("last_date"))
    .group_by(Order.customer_id)
    .subquery()
)

# Hauptabfrage
stmt = (
    select(Customer.id, Customer.name)
    .join(last_order, Customer.id == last_order.c.customer_id)
    .join(Order, (Order.customer_id == Customer.id) &
                 (Order.order_date == last_order.c.last_date))
    .where(Order.total > avg_subq)
    .order_by(Customer.name)
)
\end{lstlisting}

\begin{lstlisting}[style=javastyle, caption={Beispiel 6 -- Java (Subquery)},
  label=lst:bsp6_java]
CriteriaBuilder cb = em.getCriteriaBuilder();
CriteriaQuery<Object[]> cq = cb.createQuery(Object[].class);
Root<Customer> c = cq.from(Customer.class);
Join<Customer, Order> o = c.join("orders", JoinType.INNER);

// Scalar Subquery: Durchschnitt
Subquery<Double> avgSub = cq.subquery(Double.class);
Root<Order> so = avgSub.from(Order.class);
avgSub.select(cb.avg(so.<Double>get("total")));

cq.multiselect(c.get("id"), c.get("name"))
  .where(cb.and(
      cb.gt(o.<Double>get("total"), avgSub),
      // letzte Bestellung via correlated subquery ...
      cb.equal(o.get("orderDate"),
          cb.greatest(o.<java.time.LocalDate>get("orderDate")))
  ))
  .orderBy(cb.asc(c.get("name")));
\end{lstlisting}

\begin{lstlisting}[style=csharpstyle, caption={Beispiel 6 -- C\# (Subquery via LINQ)},
  label=lst:bsp6_cs]
var avgTotal = context.Orders.Average(o => o.Total);

var lastOrders = context.Orders
    .GroupBy(o => o.CustomerId)
    .Select(g => new {
        CustomerId = g.Key,
        LastDate   = g.Max(o => o.OrderDate)
    });

var results = from c in context.Customers
              join lo in lastOrders on c.Id equals lo.CustomerId
              join o in context.Orders
                on new { c.Id, lo.LastDate }
                equals new { Id = o.CustomerId, LastDate = o.OrderDate }
              where o.Total > avgTotal
              orderby c.Name
              select new { c.Id, c.Name };
\end{lstlisting}

\begin{successbox}
\textbf{UQTL-Normiertes SQL:}
\begin{lstlisting}[style=sqlstyle, numbers=none]
SELECT t1.id, t1.name
FROM customers AS t1
INNER JOIN orders AS t2
  ON t1.id = t2.customer_id
WHERE t2.order_date = (
    SELECT MAX(t3.order_date)
    FROM orders AS t3
    WHERE t3.customer_id = t1.id
)
  AND t2.total > (
    SELECT AVG(t4.total)
    FROM orders AS t4
)
ORDER BY t1.name ASC
\end{lstlisting}
\end{successbox}

% ─────────────────────────────────────────────────────────────────────────────
\subsection{Beispiel 7: Komplexer Multi-Join mit Aggregation}
\label{subsec:bsp7}
% ─────────────────────────────────────────────────────────────────────────────

\textbf{Aufgabe:} Top-5-Produkte nach Umsatz (Menge $\times$ Einzelpreis),
nur Produkte der Kategorie ,,Elektronik'', mit Anzahl unterschiedlicher Käufer.

\begin{lstlisting}[style=pythonstyle, caption={Beispiel 7 -- Python (Multi-Join + Aggregation)},
  label=lst:bsp7_py]
from sqlalchemy import select, func

stmt = (
    select(
        Product.id,
        Product.name,
        func.sum(OrderItem.quantity * OrderItem.unit_price).label("revenue"),
        func.count(Order.customer_id.distinct()).label("unique_buyers")
    )
    .join(OrderItem, Product.id == OrderItem.product_id)
    .join(Order, OrderItem.order_id == Order.id)
    .where(Product.category == "Elektronik")
    .group_by(Product.id, Product.name)
    .order_by(func.sum(
        OrderItem.quantity * OrderItem.unit_price).desc())
    .limit(5)
)
\end{lstlisting}

\begin{lstlisting}[style=javastyle, caption={Beispiel 7 -- Java (Multi-Join + Aggregation)},
  label=lst:bsp7_java]
CriteriaBuilder cb = em.getCriteriaBuilder();
CriteriaQuery<Object[]> cq = cb.createQuery(Object[].class);
Root<Product> p = cq.from(Product.class);
Join<Product, OrderItem> oi = p.join("orderItems", JoinType.INNER);
Join<OrderItem, Order>   o  = oi.join("order",      JoinType.INNER);

Expression<BigDecimal> revenue =
    cb.sum(cb.prod(oi.<Integer>get("quantity"),
                   oi.<BigDecimal>get("unitPrice")));
Expression<Long> buyers =
    cb.countDistinct(o.get("customerId"));

cq.multiselect(p.get("id"), p.get("name"), revenue, buyers)
  .where(cb.equal(p.get("category"), "Elektronik"))
  .groupBy(p.get("id"), p.get("name"))
  .orderBy(cb.desc(revenue));

List<Object[]> result = em.createQuery(cq)
    .setMaxResults(5)
    .getResultList();
\end{lstlisting}

\begin{lstlisting}[style=csharpstyle, caption={Beispiel 7 -- C\# (Multi-Join + Aggregation)},
  label=lst:bsp7_cs]
var results = context.Products
    .Where(p => p.Category == "Elektronik")
    .Join(context.OrderItems,
          p  => p.Id,
          oi => oi.ProductId,
          (p, oi) => new { p, oi })
    .Join(context.Orders,
          x  => x.oi.OrderId,
          o  => o.Id,
          (x, o) => new { x.p, x.oi, o })
    .GroupBy(x => new { x.p.Id, x.p.Name })
    .Select(g => new {
        g.Key.Id,
        g.Key.Name,
        Revenue      = g.Sum(x => x.oi.Quantity * x.oi.UnitPrice),
        UniqueBuyers = g.Select(x => x.o.CustomerId).Distinct().Count()
    })
    .OrderByDescending(x => x.Revenue)
    .Take(5);
\end{lstlisting}

\begin{successbox}
\textbf{UQTL-Normiertes SQL:}
\begin{lstlisting}[style=sqlstyle, numbers=none]
SELECT t1.id,
       t1.name,
       SUM(t2.quantity * t2.unit_price)   AS revenue,
       COUNT(DISTINCT t3.customer_id)     AS unique_buyers
FROM products AS t1
INNER JOIN order_items AS t2
  ON t1.id = t2.product_id
INNER JOIN orders AS t3
  ON t2.order_id = t3.id
WHERE t1.category = 'Elektronik'
GROUP BY t1.id, t1.name
ORDER BY SUM(t2.quantity * t2.unit_price) DESC
LIMIT 5
\end{lstlisting}
\end{successbox}

% ─────────────────────────────────────────────────────────────────────────────
\subsection{Beispiel 8: EXISTS-Unterabfrage}
\label{subsec:bsp8}
% ─────────────────────────────────────────────────────────────────────────────

\textbf{Aufgabe:} Alle Kunden, die mindestens eine Bestellung mit dem Status
,,shipped'' aber noch keine mit dem Status ,,delivered'' haben.

\begin{lstlisting}[style=pythonstyle, caption={Beispiel 8 -- Python (EXISTS)},
  label=lst:bsp8_py]
from sqlalchemy import select, exists

has_shipped = exists(
    select(Order.id)
    .where(Order.customer_id == Customer.id)
    .where(Order.status == "shipped")
)
has_delivered = exists(
    select(Order.id)
    .where(Order.customer_id == Customer.id)
    .where(Order.status == "delivered")
)

stmt = (
    select(Customer)
    .where(has_shipped)
    .where(~has_delivered)
    .order_by(Customer.name)
)
\end{lstlisting}

\begin{lstlisting}[style=javastyle, caption={Beispiel 8 -- Java (EXISTS)},
  label=lst:bsp8_java]
CriteriaBuilder cb = em.getCriteriaBuilder();
CriteriaQuery<Customer> cq = cb.createQuery(Customer.class);
Root<Customer> c = cq.from(Customer.class);

Subquery<Integer> shipped   = cq.subquery(Integer.class);
Root<Order> os = shipped.from(Order.class);
shipped.select(cb.literal(1))
    .where(cb.and(
        cb.equal(os.get("customerId"), c.get("id")),
        cb.equal(os.get("status"), "shipped")
    ));

Subquery<Integer> delivered = cq.subquery(Integer.class);
Root<Order> od = delivered.from(Order.class);
delivered.select(cb.literal(1))
    .where(cb.and(
        cb.equal(od.get("customerId"), c.get("id")),
        cb.equal(od.get("status"), "delivered")
    ));

cq.select(c)
  .where(cb.and(
      cb.exists(shipped),
      cb.not(cb.exists(delivered))
  ))
  .orderBy(cb.asc(c.get("name")));
\end{lstlisting}

\begin{lstlisting}[style=csharpstyle, caption={Beispiel 8 -- C\# (EXISTS via Any)},
  label=lst:bsp8_cs]
var results = context.Customers
    .Where(c =>
        context.Orders.Any(o => o.CustomerId == c.Id
                             && o.Status == "shipped") &&
       !context.Orders.Any(o => o.CustomerId == c.Id
                             && o.Status == "delivered"))
    .OrderBy(c => c.Name);
\end{lstlisting}

\begin{successbox}
\textbf{UQTL-Normiertes SQL:}
\begin{lstlisting}[style=sqlstyle, numbers=none]
SELECT t1.id, t1.name, t1.email, t1.country,
       t1.active, t1.created_at, t1.tier
FROM customers AS t1
WHERE EXISTS (
    SELECT 1
    FROM orders AS t2
    WHERE t2.customer_id = t1.id
      AND t2.status = 'shipped'
)
  AND NOT EXISTS (
    SELECT 1
    FROM orders AS t3
    WHERE t3.customer_id = t1.id
      AND t3.status = 'delivered'
)
ORDER BY t1.name ASC
\end{lstlisting}
\end{successbox}

% ─────────────────────────────────────────────────────────────────────────────
\subsection{Beispiel 9: Common Table Expressions (WITH)}
\label{subsec:bsp9}
% ─────────────────────────────────────────────────────────────────────────────

\textbf{Aufgabe:} Kunden mit überdurchschnittlichem Umsatz (Top-Tier-Analyse),
unter Verwendung von CTEs für Lesbarkeit.

\begin{lstlisting}[style=pythonstyle, caption={Beispiel 9 -- Python (CTE)},
  label=lst:bsp9_py]
from sqlalchemy import select, func

# CTE 1: Gesamtumsatz je Kunde
customer_revenue = (
    select(
        Order.customer_id,
        func.sum(Order.total).label("total_revenue")
    )
    .group_by(Order.customer_id)
    .cte("customer_revenue")
)

# CTE 2: Durchschnittsumsatz
avg_revenue = (
    select(func.avg(customer_revenue.c.total_revenue).label("avg_rev"))
    .cte("avg_revenue")
)

# Hauptabfrage
stmt = (
    select(Customer.name, Customer.email,
           customer_revenue.c.total_revenue)
    .join(customer_revenue,
          Customer.id == customer_revenue.c.customer_id)
    .where(customer_revenue.c.total_revenue >
           select(avg_revenue.c.avg_rev).scalar_subquery())
    .order_by(customer_revenue.c.total_revenue.desc())
)
\end{lstlisting}

\begin{lstlisting}[style=csharpstyle, caption={Beispiel 9 -- C\# (CTE-Aequivalent)},
  label=lst:bsp9_cs]
// EF Core: CTEs werden als verkettete LINQ-Ausdruecke modelliert
var customerRevenue = context.Orders
    .GroupBy(o => o.CustomerId)
    .Select(g => new {
        CustomerId    = g.Key,
        TotalRevenue  = g.Sum(o => o.Total)
    });

var avgRevenue = customerRevenue.Average(x => x.TotalRevenue);

var results = context.Customers
    .Join(customerRevenue,
          c  => c.Id,
          cr => cr.CustomerId,
          (c, cr) => new { c, cr })
    .Where(x => x.cr.TotalRevenue > avgRevenue)
    .OrderByDescending(x => x.cr.TotalRevenue)
    .Select(x => new {
        x.c.Name,
        x.c.Email,
        x.cr.TotalRevenue
    });
\end{lstlisting}

\begin{successbox}
\textbf{UQTL-Normiertes SQL (mit CTE):}
\begin{lstlisting}[style=sqlstyle, numbers=none]
WITH customer_revenue AS (
    SELECT t2.customer_id,
           SUM(t2.total) AS total_revenue
    FROM orders AS t2
    GROUP BY t2.customer_id
),
avg_revenue AS (
    SELECT AVG(cr.total_revenue) AS avg_rev
    FROM customer_revenue AS cr
)
SELECT t1.name, t1.email, cr.total_revenue
FROM customers AS t1
INNER JOIN customer_revenue AS cr
  ON t1.id = cr.customer_id
WHERE cr.total_revenue > (SELECT avg_rev FROM avg_revenue)
ORDER BY cr.total_revenue DESC
\end{lstlisting}
\end{successbox}

% ─────────────────────────────────────────────────────────────────────────────
\subsection{Beispiel 10: Vollständige Geschäftslogik-Abfrage}
\label{subsec:bsp10}
% ─────────────────────────────────────────────────────────────────────────────

\textbf{Aufgabe:} Monatlicher Umsatzbericht für das Jahr 2024: Monat, Gesamtumsatz,
Anzahl Bestellungen, Anzahl eindeutiger Kunden, Vergleich zum Vormonat
(via LAG-Fensterfunktion).

\begin{lstlisting}[style=pythonstyle, caption={Beispiel 10 -- Python (Business Report)},
  label=lst:bsp10_py]
from sqlalchemy import select, func, extract

monthly = (
    select(
        extract("year",  Order.order_date).label("year"),
        extract("month", Order.order_date).label("month"),
        func.sum(Order.total).label("revenue"),
        func.count(Order.id).label("order_count"),
        func.count(Order.customer_id.distinct()).label("unique_customers")
    )
    .where(extract("year", Order.order_date) == 2024)
    .where(Order.status != "pending")
    .group_by(
        extract("year",  Order.order_date),
        extract("month", Order.order_date)
    )
    .order_by(
        extract("year",  Order.order_date),
        extract("month", Order.order_date)
    )
    .subquery()
)

prev_rev = func.lag(monthly.c.revenue, 1).over(
    order_by=[monthly.c.year, monthly.c.month]
).label("prev_month_revenue")

stmt = select(monthly, prev_rev)
\end{lstlisting}

\begin{lstlisting}[style=csharpstyle, caption={Beispiel 10 -- C\# (Business Report)},
  label=lst:bsp10_cs]
var monthly = context.Orders
    .Where(o => o.OrderDate.Year == 2024
             && o.Status != "pending")
    .GroupBy(o => new {
        Year  = o.OrderDate.Year,
        Month = o.OrderDate.Month
    })
    .Select(g => new {
        g.Key.Year,
        g.Key.Month,
        Revenue         = g.Sum(o => o.Total),
        OrderCount      = g.Count(),
        UniqueCustomers = g.Select(o => o.CustomerId).Distinct().Count()
    })
    .OrderBy(x => x.Year)
    .ThenBy(x => x.Month);

// LAG-Fenster: EF Core 8+
var withLag = monthly
    .Select(m => new {
        m.Year, m.Month, m.Revenue, m.OrderCount, m.UniqueCustomers,
        PrevMonthRevenue = EF.Functions.Lag(m.Revenue, 1)
            .Over()
            .OrderBy(m.Year).ThenBy(m.Month)
    });
\end{lstlisting}

\begin{successbox}
\textbf{UQTL-Normiertes SQL:}
\begin{lstlisting}[style=sqlstyle, numbers=none]
SELECT sub.year,
       sub.month,
       sub.revenue,
       sub.order_count,
       sub.unique_customers,
       LAG(sub.revenue, 1) OVER (
           ORDER BY sub.year ASC, sub.month ASC
       ) AS prev_month_revenue
FROM (
    SELECT EXTRACT(YEAR  FROM t1.order_date) AS year,
           EXTRACT(MONTH FROM t1.order_date) AS month,
           SUM(t1.total)                      AS revenue,
           COUNT(t1.id)                       AS order_count,
           COUNT(DISTINCT t1.customer_id)     AS unique_customers
    FROM orders AS t1
    WHERE EXTRACT(YEAR FROM t1.order_date) = 2024
      AND t1.status != 'pending'
    GROUP BY EXTRACT(YEAR  FROM t1.order_date),
             EXTRACT(MONTH FROM t1.order_date)
    ORDER BY EXTRACT(YEAR  FROM t1.order_date) ASC,
             EXTRACT(MONTH FROM t1.order_date) ASC
) AS sub
\end{lstlisting}
\end{successbox}

\newpage
% ═══════════════════════════════════════════════════════════════════════════════
\section{Konformitätsstufen}
\label{sec:konformitaet}
% ═══════════════════════════════════════════════════════════════════════════════

Der UQTL-Standard definiert vier Konformitätsstufen (Conformance Levels),
die eine schrittweise Implementierung ermöglichen. Höhere Stufen schließen
alle Anforderungen niedrigerer Stufen ein.

\begin{figure}[H]
\centering
\begin{tikzpicture}[
  level/.style={draw, rounded corners=4pt, minimum width=11cm,
    minimum height=1.1cm, align=left, font=\small,
    inner sep=8pt},
  cl1/.style={level, fill=green!20},
  cl2/.style={level, fill=blue!15},
  cl3/.style={level, fill=orange!20},
  cl4/.style={level, fill=red!15},
]
  \node[cl1] (l1) at (0,0)   {\textbf{CL-1 (Basis):}~~~Filter, Projektion, Sortierung, LIMIT/OFFSET};
  \node[cl2] (l2) at (0,1.3) {\textbf{CL-2 (Standard):}~~~CL-1 + Aggregation, GROUP BY, HAVING, INNER JOIN};
  \node[cl3] (l3) at (0,2.6) {\textbf{CL-3 (Erweitert):}~~~CL-2 + alle JOIN-Typen, Unterabfragen, EXISTS, UNION};
  \node[cl4] (l4) at (0,3.9) {\textbf{CL-4 (Vollständig):}~~~CL-3 + Fensterfunktionen, CTEs, Mengenoperationen};
\end{tikzpicture}
\caption{UQTL-Konformitätspyramide (CL-1 bis CL-4)}
\label{fig:konformitaet}
\end{figure}

\subsection{CL-1: Basiskonformität}
\label{subsec:cl1}

Ein CL-1-konformes System muss die folgenden Operatoren vollständig und
äquivalenzerhaltend unterstützen:

\begin{itemize}[leftmargin=1.5cm]
  \item \textbf{SCAN:} Zugriff auf einzelne Tabellen mit optionalem Alias
  \item \textbf{FILTER (WHERE):} Alle Vergleichsoperatoren ($=$, $\neq$, $<$, $\leq$, $>$, $\geq$),
    \texttt{IN}, \texttt{BETWEEN}, \texttt{LIKE}, \texttt{IS NULL}, \texttt{AND}, \texttt{OR}, \texttt{NOT}
  \item \textbf{PROJECT:} Spaltenselektion, Ausdrücke, \texttt{DISTINCT}
  \item \textbf{SORT:} \texttt{ORDER BY} ASC/DESC, mehrstufig
  \item \textbf{LIMIT:} \texttt{LIMIT n OFFSET m}
\end{itemize}

\subsection{CL-2: Standardkonformität}
\label{subsec:cl2}

Zusätzlich zu CL-1 müssen CL-2-konforme Systeme unterstützen:

\begin{itemize}[leftmargin=1.5cm]
  \item \textbf{GROUP\_AGG:} \texttt{GROUP BY}, alle Standardaggregationen (\texttt{COUNT}, \texttt{SUM}, \texttt{AVG}, \texttt{MIN}, \texttt{MAX})
  \item \textbf{FILTER (HAVING):} Aggregatprädikate nach \texttt{GROUP BY}
  \item \textbf{JOIN (INNER):} \texttt{INNER JOIN ... ON}, \texttt{NATURAL JOIN}
  \item \textbf{Arithmetische Ausdrücke} in Projektionen und Prädikaten
\end{itemize}

\subsection{CL-3: Erweiterte Konformität}
\label{subsec:cl3}

Zusätzlich zu CL-2:

\begin{itemize}[leftmargin=1.5cm]
  \item \textbf{JOIN (alle Typen):} \texttt{LEFT/RIGHT/FULL OUTER JOIN}, \texttt{CROSS JOIN}
  \item \textbf{SUBQUERY:} Korrelierte und nicht-korrelierte Unterabfragen, Scalar Subqueries
  \item \textbf{EXISTS / NOT EXISTS}
  \item \textbf{SET\_OP:} \texttt{UNION}, \texttt{UNION ALL}
  \item \textbf{Skalare Funktionen:} Zeichenkettenfunktionen, Datumsfunktionen, mathematische Funktionen
\end{itemize}

\subsection{CL-4: Vollständige Konformität}
\label{subsec:cl4}

Zusätzlich zu CL-3:

\begin{itemize}[leftmargin=1.5cm]
  \item \textbf{WINDOW:} Alle ANSI-Fensterfunktionen (\texttt{ROW\_NUMBER}, \texttt{RANK}, \texttt{DENSE\_RANK}, \texttt{LAG}, \texttt{LEAD}, \texttt{FIRST\_VALUE}, \texttt{LAST\_VALUE}, \texttt{NTH\_VALUE}, gleitende Aggregate)
  \item \textbf{WITH (CTEs):} Einfache und rekursive CTEs
  \item \textbf{SET\_OP vollständig:} \texttt{INTERSECT}, \texttt{EXCEPT}
  \item \textbf{CASE-Ausdrücke}
  \item \textbf{Dialekt-Erweiterungen} (als optionale Zertifizierung)
\end{itemize}

\begin{table}[H]
\centering
\caption{Konformitätsstufen-Übersicht}
\label{tab:konformitaet}
\renewcommand{\arraystretch}{1.3}
\begin{tabularx}{\textwidth}{lXXXX}
\toprule
\textbf{Feature} & \textbf{CL-1} & \textbf{CL-2} & \textbf{CL-3} & \textbf{CL-4} \\
\midrule
Filter (WHERE) & \checkmark & \checkmark & \checkmark & \checkmark \\
Projektion & \checkmark & \checkmark & \checkmark & \checkmark \\
Sortierung & \checkmark & \checkmark & \checkmark & \checkmark \\
LIMIT/OFFSET & \checkmark & \checkmark & \checkmark & \checkmark \\
GROUP BY / Aggregation & & \checkmark & \checkmark & \checkmark \\
HAVING & & \checkmark & \checkmark & \checkmark \\
INNER JOIN & & \checkmark & \checkmark & \checkmark \\
OUTER JOINs & & & \checkmark & \checkmark \\
Unterabfragen & & & \checkmark & \checkmark \\
EXISTS & & & \checkmark & \checkmark \\
UNION / UNION ALL & & & \checkmark & \checkmark \\
Fensterfunktionen & & & & \checkmark \\
CTEs (WITH) & & & & \checkmark \\
INTERSECT / EXCEPT & & & & \checkmark \\
CASE-Ausdrücke & & & & \checkmark \\
\bottomrule
\end{tabularx}
\end{table}

% ═══════════════════════════════════════════════════════════════════════════════
\section{Referenzimplementierung}
\label{sec:implementierung}
% ═══════════════════════════════════════════════════════════════════════════════

\subsection{Architektur einer UQTL-konformen Implementierung}
\label{subsec:impl_arch}

Eine UQTL-konforme Implementierung besteht aus drei Schichten:

\begin{enumerate}[leftmargin=1.5cm]
  \item \textbf{Frontend-Parser:} Sprach\-spezifische Analyse und
    IQT-Aufbau (je ein Parser pro Quellsprache)
  \item \textbf{IQT-Normalisierer:} Anwendung der Normalisierungsregeln
    (\ref{subsec:iqt_norm})
  \item \textbf{SQL-Codegenerator:} IQT-zu-SQL-Übersetzung nach Zieldialekt
\end{enumerate}

\subsection{Pseudocode des IQT-Builders (Python-Frontend)}
\label{subsec:impl_py}

\begin{lstlisting}[style=pythonstyle, caption={Python-Frontend: IQT-Builder (vereinfacht)},
  label=lst:impl_py]
class PythonIQTBuilder:
    """Baut einen IQT-Baum aus einem SQLAlchemy Select-Statement."""

    def build(self, stmt: Select) -> IQTNode:
        node = self._visit_froms(stmt)          # SCAN / JOIN
        node = self._visit_where(stmt, node)    # FILTER WHERE
        node = self._visit_group_having(stmt, node)  # GROUP_AGG + HAVING
        node = self._visit_select(stmt, node)   # PROJECT
        node = self._visit_order(stmt, node)    # SORT
        node = self._visit_limit(stmt, node)    # LIMIT
        return node

    def _visit_where(self, stmt, source):
        if stmt.whereclause is not None:
            pred = self._translate_predicate(stmt.whereclause)
            return FilterNode(source=source, predicate=pred,
                              position=FilterPosition.WHERE)
        return source

    def _translate_predicate(self, clause) -> IQTExpr:
        if isinstance(clause, BinaryExpression):
            left  = self._translate_expr(clause.left)
            right = self._translate_expr(clause.right)
            op    = self._map_operator(clause.operator)
            return CompareExpr(left=left, op=op, right=right)
        elif isinstance(clause, BooleanClauseList):
            children = [self._translate_predicate(c) for c in clause.clauses]
            logic = LogicOp.AND if clause.operator == operators.and_ \
                    else LogicOp.OR
            return LogicExpr(op=logic, operands=children)
        # ... weitere Faelle
\end{lstlisting}

\subsection{Pseudocode des SQL-Codegenerators}
\label{subsec:impl_codegen}

\begin{lstlisting}[style=pythonstyle, caption={SQL-Codegenerator: IQT nach ANSI SQL},
  label=lst:impl_codegen]
class AnsiSqlCodegen:
    """Generiert ANSI SQL:2016 aus einem normalisierten IQT."""

    def generate(self, node: IQTNode) -> str:
        parts = self._collect(node)
        return self._assemble(parts)

    def _collect(self, node, parts=None) -> SqlParts:
        if parts is None:
            parts = SqlParts()

        if isinstance(node, ScanNode):
            parts.from_clause = f"{node.table} AS {node.alias}"

        elif isinstance(node, FilterNode):
            self._collect(node.source, parts)
            pred_sql = self._gen_expr(node.predicate)
            if node.position == FilterPosition.WHERE:
                parts.where_clauses.append(pred_sql)
            else:
                parts.having_clauses.append(pred_sql)

        elif isinstance(node, JoinNode):
            self._collect(node.left, parts)
            right_sql  = f"{node.right.table} AS {node.right.alias}"
            cond_sql   = self._gen_expr(node.condition.predicate)
            join_kw    = JOIN_KEYWORDS[node.join_type]
            parts.joins.append(f"{join_kw} {right_sql} ON {cond_sql}")

        elif isinstance(node, GroupAggNode):
            self._collect(node.source, parts)
            parts.group_by = ", ".join(
                self._gen_expr(e) for e in node.group_by)
            parts.select_items.extend(
                f"{self._gen_agg(a)} AS {a.alias}" for a in node.aggregates)

        elif isinstance(node, ProjectNode):
            self._collect(node.source, parts)
            parts.distinct = node.distinct
            parts.select_items = [
                self._gen_proj(item) for item in node.columns]

        elif isinstance(node, SortNode):
            self._collect(node.source, parts)
            parts.order_by = ", ".join(
                self._gen_sort_item(s) for s in node.items)

        elif isinstance(node, LimitNode):
            self._collect(node.source, parts)
            if node.count is not None:
                parts.limit  = str(self._gen_expr(node.count))
            if node.offset is not None:
                parts.offset = str(self._gen_expr(node.offset))

        return parts

    def _assemble(self, p: SqlParts) -> str:
        distinct = "DISTINCT " if p.distinct else ""
        sel = ", ".join(p.select_items) if p.select_items else "*"
        sql = f"SELECT {distinct}{sel}\nFROM {p.from_clause}"
        if p.joins:
            sql += "\n" + "\n".join(p.joins)
        if p.where_clauses:
            sql += "\nWHERE " + "\n  AND ".join(p.where_clauses)
        if p.group_by:
            sql += f"\nGROUP BY {p.group_by}"
        if p.having_clauses:
            sql += "\nHAVING " + "\n  AND ".join(p.having_clauses)
        if p.order_by:
            sql += f"\nORDER BY {p.order_by}"
        if p.limit:
            sql += f"\nLIMIT {p.limit}"
        if p.offset:
            sql += f"\nOFFSET {p.offset}"
        return sql
\end{lstlisting}

\subsection{Konformitätstests}
\label{subsec:impl_tests}

Eine UQTL-Referenzimplementierung muss eine Test-Suite bestehen, die für
jede Konformitätsstufe die semantische Äquivalenz der Übersetzungen prüft:

\begin{lstlisting}[style=pythonstyle, caption={Konformitätstest-Schablone},
  label=lst:impl_test]
class UQTLConformanceTest:

    def assert_equivalent(self, py_stmt, java_iqt, cs_iqt, expected_sql):
        """Prueft semantische Aequivalenz aller drei Uebersetzungen."""
        py_sql   = PythonIQTBuilder().build(py_stmt)
                   |> IQTNormalizer().normalize()
                   |> AnsiSqlCodegen().generate()
        java_sql = java_iqt
                   |> IQTNormalizer().normalize()
                   |> AnsiSqlCodegen().generate()
        cs_sql   = cs_iqt
                   |> IQTNormalizer().normalize()
                   |> AnsiSqlCodegen().generate()

        assert py_sql   == expected_sql, f"Python divergiert:\n{py_sql}"
        assert java_sql == expected_sql, f"Java divergiert:\n{java_sql}"
        assert cs_sql   == expected_sql, f"C# divergiert:\n{cs_sql}"
\end{lstlisting}

\newpage
% sec7_implementierung.tex (leer, wird durch sec6 abgedeckt)

\newpage
% ═══════════════════════════════════════════════════════════════════════════════
\section{Diskussion}
\label{sec:diskussion}
% ═══════════════════════════════════════════════════════════════════════════════

\subsection{Stärken des UQTL-Standards}
\label{subsec:staerken}

Der UQTL-Standard bietet gegenüber dem Status quo mehrere fundamentale
Vorteile:

\paragraph{Verifikationsgrundlage}
Durch das definierte Normalformat können Query-Äquivalenz\-prüfungen vollständig
automatisiert werden. Dies ermöglicht:
\begin{itemize}[leftmargin=1.5cm]
  \item Automatisierte Regressionstests bei Frameworkmigrationen
  \item Statische Analyse von Abfrage-Äquivalenz in CI/CD-Pipelines
  \item Formale Verifikation kritischer Datenzugriffspfade
\end{itemize}

\paragraph{Reduzierung kognitiver Last}
Entwickler, die in mehrsprachigen Teams arbeiten, können Abfragen in ihrer
bevorzugten Sprache formulieren und sind dennoch garantiert, dass semantisch
gleiche Abfragen zu identischen SQL-Anweisungen führen. Code-Reviews werden
deutlich vereinfacht, da der Vergleich auf der Ebene des normalisierten SQL
stattfinden kann.

\paragraph{Langzeitstabilität}
Frameworks evolvieren: SQLAlchemy hat mit Version 2.0 eine neue API eingeführt,
Entity Framework hat mehrere inkompatible Generationen durchlaufen. Ein
stabiler Standard abstrahiert von diesen Evolutionen.

\subsection{Grenzen und Limitierungen}
\label{subsec:grenzen}

\paragraph{Semantisch nicht äquivalente Quellkonstrukte}
Nicht alle Quellsprachkonstrukte haben eindeutige SQL-Entsprechungen.
Beispielsweise kann eine Python-Lambda-Funktion in einem \texttt{filter()}-Ausdruck
beliebige Python-Logik enthalten, die kein SQL-Äquivalent hat. UQTL definiert
für solche Fälle einen \emph{Fallback-Mechanismus}: Der betreffende Ausdruck
wird als \texttt{OPAQUE}-Knoten im IQT markiert und löst einen Normalisierungsfehler
aus.

\paragraph{Dialektspezifische Funktionen}
Bestimmte Datenbankfunktionen (z.\,B.\ PostgreSQL-spezifisches \texttt{jsonb\_agg()},
Oracle-spezifisches \texttt{CONNECT BY}) sind außerhalb von Erweiterungen nicht
im Standard enthalten. Implementierungen, die diese benötigen, müssen einen
Dialekt-Erweiterungsmechanismus implementieren.

\paragraph{Performanz-Implikationen}
Die Normalisierung erzeugt in einigen Fällen SQL-Ausdrücke, die semantisch
korrekt, aber nicht optimal für alle Datenbankplaner sind. Implementierungen
sind ermutigt, nach der Normalisierung einen optionalen \emph{Hint-Layer}
anzuwenden.

\paragraph{Stream-API vs. Datenbankabfragen}
Java Streams operieren in-memory und unterstützen Operationen (z.\,B.\
\texttt{peek()}, \texttt{forEach()}, zustandsbehaftete Lambdas), die keine
SQL-Entsprechung haben. UQTL richtet sich an den Subset der Stream-Operationen,
der eine relationale Semantik trägt.

\subsection{Erweiterungen und zukünftige Arbeit}
\label{subsec:erweiterungen}

Folgende Erweiterungen sind für zukünftige Versionen des Standards geplant:

\begin{enumerate}[leftmargin=1.5cm]
  \item \textbf{UQTL-E1 (NoSQL-Dialekte):} Erweiterung um MongoDB-Aggregations\-pipeline
    und Elasticsearch-DSL als zusätzliche Quellsprachen
  \item \textbf{UQTL-E2 (TypeScript):} Unterstützung von TypeScript/Prisma als
    vierte Quellsprache
  \item \textbf{UQTL-E3 (Streaming SQL):} Integration von Apache Flink SQL und
    Apache Spark SQL als Zieldialekte
  \item \textbf{UQTL-E4 (Schema-Awareness):} Integration von Schemavalidierung
    und Typ-Prüfung in den Normalisierungsprozess
  \item \textbf{UQTL-E5 (ML-gestützte Optimierung):} Automatische Query-Optimierung
    basierend auf historischen Ausführungsstatistiken
\end{enumerate}

% ═══════════════════════════════════════════════════════════════════════════════
\section{Zusammenfassung}
\label{sec:zusammenfassung}
% ═══════════════════════════════════════════════════════════════════════════════

Die vorliegende Arbeit hat den \textbf{Unified Query Translation Layer (UQTL)}
vollständig spezifiziert: einen offenen Standard zur garantierten semantischen
Äquivalenz von Query-Übersetzungen aus Python, Java (Stream API / JPA) und
C\#/.NET (LINQ) nach SQL.

\paragraph{Kernbeiträge}
Die wesentlichen Beiträge des Standards lassen sich wie folgt zusammenfassen:

\begin{enumerate}[leftmargin=1.5cm]
  \item \textbf{Formale Grundlage:} Die Definition des Intermediate Query Tree
    (IQT) als typisierter DAG mit vollständigen Knotenspezifikationen für alle
    11~Knotentypen (SCAN, FILTER, PROJECT, JOIN, GROUP\_AGG, SORT, LIMIT,
    WINDOW, SUBQUERY, SET\_OP, WITH).

  \item \textbf{Vollständige Übersetzungsregeln:} Für alle 9~Hauptoperatoren
    (Filter, Projektion, Sortierung, Limit, Aggregation, Joins, Fensterfunktionen,
    Unterabfragen, Mengenoperationen) wurden vollständige, formale
    Ãœbersetzungsregeln von allen drei Quellsprachen in den IQT und vom IQT
    nach ANSI SQL:2016 definiert.

  \item \textbf{10~Praxisbeispiele:} Anhand eines einheitlichen E-Commerce-Daten\-modells
    wurden 10~vollständige Beispiele demonstriert, die von einfacher Filterung bis
    zu komplexen Geschäftsberichten mit CTEs und Fensterfunktionen reichen.

  \item \textbf{Konformitätsstufen:} Vier Stufen (CL-1 bis CL-4) erlauben
    einen schrittweisen Aufbau konformer Implementierungen.

  \item \textbf{Referenzimplementierungsrichtlinien:} Architektur und
    Pseudocode einer konformen Implementierung wurden bereitgestellt.
\end{enumerate}

\paragraph{Bedeutung für die Praxis}

Der UQTL-Standard schließt eine fundamentale Lücke in der polyglotten
Softwareentwicklung: Erstmals können Teams aus Python-, Java- und
C\#-Entwicklern formal garantieren, dass ihre Datenbankzugriffe semantisch
äquivalent sind -- unabhängig von Sprache, Framework-Version oder persönlichem
Codingstil.

\begin{successbox}
\textbf{Leitprinzip des Standards:}

\emph{Dieselbe Datenabfrage, formuliert in unterschiedlichen Sprachen,
erzeugt unter UQTL stets dieselbe SQL-Anweisung. Die Sprache ist ein
Implementierungsdetail -- die Semantik ist der Standard.}
\end{successbox}


\newpage
% Platzhalter -- Inhalt in sec8_diskussion.tex integriert

\newpage
% ═══════════════════════════════════════════════════════════════════════════════
\appendix
\section{Vollständige IQT-Grammatik (EBNF)}
\label{app:grammatik}
% ═══════════════════════════════════════════════════════════════════════════════

Die folgende EBNF-Grammatik definiert die vollständige Syntax des
Intermediate Query Tree.

\begin{lstlisting}[style=sqlstyle, caption={IQT-Grammatik in EBNF},
  label=lst:iqt_ebnf, numbers=left]
IQT        ::= QueryNode

QueryNode  ::= ScanNode
             | FilterNode
             | ProjectNode
             | JoinNode
             | GroupAggNode
             | SortNode
             | LimitNode
             | WindowNode
             | SubqueryNode
             | SetOpNode
             | WithNode

ScanNode       ::= 'SCAN' '{' 'table' ':' QualName
                              ('alias' ':' Ident)?
                              ('hint' ':' IndexHint)?
                   '}'

FilterNode     ::= 'FILTER' '{' 'source'    ':' QueryNode
                                'predicate' ':' Expr
                                'position'  ':' ('WHERE' | 'HAVING')
                   '}'

ProjectNode    ::= 'PROJECT' '{' 'source'   ':' QueryNode
                                 'columns'  ':' ProjItem+
                                 'distinct' ':' Bool
                   '}'

ProjItem       ::= 'Star' | 'Expr' '{' 'expr' ':' Expr ('alias' ':' Ident)? '}'

JoinNode       ::= 'JOIN' '{' 'left'      ':' QueryNode
                              'right'     ':' QueryNode
                              'type'      ':' JoinType
                              'condition' ':' JoinCond
                   '}'

JoinType       ::= 'INNER' | 'LEFT_OUTER' | 'RIGHT_OUTER'
                 | 'FULL_OUTER' | 'CROSS'

JoinCond       ::= 'ON'    '{' 'predicate' ':' Expr '}'
                 | 'USING' '{' 'columns'   ':' Ident+ '}'
                 | 'NATURAL'

GroupAggNode   ::= 'GROUP_AGG' '{' 'source'     ':' QueryNode
                                   'groupBy'    ':' Expr*
                                   'aggregates' ':' AggItem+
                                   ('having'    ':' Expr)?
                   '}'

AggItem        ::= '{' 'func'     ':' AggFunc
                       'arg'      ':' (Expr | 'STAR')
                       'distinct' ':' Bool
                       ('alias'   ':' Ident)?
                   '}'

AggFunc        ::= 'COUNT' | 'SUM' | 'AVG' | 'MIN' | 'MAX'
                 | 'STDDEV' | 'VAR' | 'STRING_AGG' | 'ARRAY_AGG'

SortNode       ::= 'SORT' '{' 'source' ':' QueryNode
                              'items'  ':' SortItem+
                   '}'

SortItem       ::= '{' 'expr'      ':' Expr
                       'direction' ':' ('ASC' | 'DESC')
                       'nulls'     ':' ('NULLS_FIRST' | 'NULLS_LAST' | 'DEFAULT')
                   '}'

LimitNode      ::= 'LIMIT' '{' 'source' ':' QueryNode
                               'count'  ':' (Expr | 'NULL')
                               'offset' ':' (Expr | 'NULL')
                   '}'

WindowNode     ::= 'WINDOW' '{' 'source'      ':' QueryNode
                                'definitions' ':' WindowDef+
                   '}'

WindowDef      ::= '{' 'alias'       ':' Ident
                       ('partitionBy' ':' Expr*)?
                       ('orderBy'     ':' SortItem*)?
                       ('frame'       ':' WindowFrame)?
                       'func'         ':' WindowFunc
                       'resultAlias'  ':' Ident
                   '}'

WindowFrame    ::= '{' 'unit'  ':' ('ROWS' | 'RANGE')
                       'start' ':' FrameBound
                       'end'   ':' FrameBound
                   '}'

FrameBound     ::= 'UNBOUNDED_PRECEDING' | 'CURRENT_ROW'
                 | 'UNBOUNDED_FOLLOWING'
                 | 'PRECEDING' '{' Expr '}'
                 | 'FOLLOWING' '{' Expr '}'

SetOpNode      ::= 'SET_OP' '{' 'left'  ':' QueryNode
                                'right' ':' QueryNode
                                'op'    ':' ('UNION' | 'UNION_ALL'
                                            | 'INTERSECT' | 'EXCEPT')
                   '}'

WithNode       ::= 'WITH' '{' 'ctes'  ':' CTEDef+
                              'query' ':' QueryNode
                   '}'

CTEDef         ::= '{' 'name'      ':' Ident
                       'recursive' ':' Bool
                       'query'     ':' QueryNode
                   '}'

Expr           ::= CompareExpr | LogicExpr | InExpr | BetweenExpr
                 | LikeExpr | NullExpr | ExistsExpr | CaseExpr
                 | Literal | ColumnRef | FuncCall | ArithExpr | SubqueryExpr

CompareExpr    ::= '{' 'left' ':' Expr 'op' ':' CmpOp 'right' ':' Expr '}'
CmpOp          ::= 'EQ' | 'NEQ' | 'LT' | 'GT' | 'LTE' | 'GTE'
LogicExpr      ::= '{' 'op' ':' ('AND'|'OR'|'NOT') 'operands' ':' Expr+ '}'
InExpr         ::= '{' 'value' ':' Expr 'list' ':' (Expr+ | SubqueryNode)
                       'negated' ':' Bool '}'
BetweenExpr    ::= '{' 'value' ':' Expr 'lo' ':' Expr 'hi' ':' Expr
                       'negated' ':' Bool '}'
LikeExpr       ::= '{' 'value' ':' Expr 'pattern' ':' StringLit
                       ('escape' ':' Char)? 'negated' ':' Bool '}'
NullExpr       ::= '{' 'value' ':' Expr 'negated' ':' Bool '}'
ExistsExpr     ::= '{' 'subquery' ':' SubqueryNode 'negated' ':' Bool '}'
\end{lstlisting}

% ─────────────────────────────────────────────────────────────────────────────
\section{Operatormapping-Referenztabelle}
\label{app:mapping}
% ─────────────────────────────────────────────────────────────────────────────

\begin{landscape}
\scriptsize
\begin{longtable}{p{4.5cm}p{5.5cm}p{5.5cm}p{5.5cm}}
\caption{Vollständige Operatormapping-Referenz (IQT $\leftrightarrow$ Quellsprachen)}
\label{tab:full_mapping} \\
\toprule
\textbf{IQT-Konstrukt} & \textbf{Python (SA 2.0)} & \textbf{Java (JPA/QueryDSL)} & \textbf{C\# (LINQ)} \\
\midrule
\endfirsthead
\multicolumn{4}{c}{\small\itshape Fortsetzung von Tabelle~\ref{tab:full_mapping}} \\
\toprule
\textbf{IQT-Konstrukt} & \textbf{Python} & \textbf{Java} & \textbf{C\#} \\
\midrule
\endhead
\midrule
\multicolumn{4}{r}{\small\itshape Fortsetzung auf nächster Seite} \\
\endfoot
\bottomrule
\endlastfoot

SCAN (table)
  & \texttt{select(Model)}
  & \texttt{cq.from(Model.class)}
  & \texttt{context.Models} \\[4pt]

SCAN (alias)
  & \texttt{aliased(Model, name="a")}
  & \texttt{Root<M> a = cq.from(M.class)}
  & \texttt{-- (implizit)} \\[4pt]

FILTER (=)
  & \texttt{.where(M.col == v)}
  & \texttt{.where(cb.equal(r.get("col"),v))}
  & \texttt{.Where(x => x.Col == v)} \\[4pt]

FILTER (!=)
  & \texttt{.where(M.col != v)}
  & \texttt{.where(cb.notEqual(...))}
  & \texttt{.Where(x => x.Col != v)} \\[4pt]

FILTER (LIKE)
  & \texttt{.where(M.col.like("\%x\%"))}
  & \texttt{.where(cb.like(r.get("col"),"\%x\%"))}
  & \texttt{.Where(x => EF.Functions.Like(x.Col,"\%x\%"))} \\[4pt]

FILTER (IS NULL)
  & \texttt{.where(M.col == None)}
  & \texttt{.where(cb.isNull(r.get("col")))}
  & \texttt{.Where(x => x.Col == null)} \\[4pt]

FILTER (IN list)
  & \texttt{.where(M.col.in\_(list))}
  & \texttt{.where(r.get("col").in(list))}
  & \texttt{.Where(x => list.Contains(x.Col))} \\[4pt]

FILTER (IN subq)
  & \texttt{.where(M.col.in\_(subq))}
  & \texttt{.where(r.get("col").in(subq))}
  & \texttt{.Where(x => subq.Contains(x.Col))} \\[4pt]

FILTER (BETWEEN)
  & \texttt{.where(M.col.between(a,b))}
  & \texttt{.where(cb.between(r.get("col"),a,b))}
  & \texttt{.Where(x => x.Col >= a \&\& x.Col <= b)} \\[4pt]

FILTER (EXISTS)
  & \texttt{.where(exists(subq))}
  & \texttt{.where(cb.exists(subq))}
  & \texttt{.Where(x => subq.Any(...))} \\[4pt]

PROJECT (cols)
  & \texttt{select(M.a, M.b)}
  & \texttt{cq.multiselect(r.get("a"),r.get("b"))}
  & \texttt{.Select(x => new\{x.A,x.B\})} \\[4pt]

PROJECT (expr)
  & \texttt{select(func.upper(M.name))}
  & \texttt{cq.select(cb.upper(r.get("name")))}
  & \texttt{.Select(x => x.Name.ToUpper())} \\[4pt]

PROJECT (DISTINCT)
  & \texttt{select(M.col).distinct()}
  & \texttt{cq.select(r.get("col")).distinct(true)}
  & \texttt{.Select(x => x.Col).Distinct()} \\[4pt]

SORT (ASC)
  & \texttt{.order\_by(M.col)}
  & \texttt{.orderBy(cb.asc(r.get("col")))}
  & \texttt{.OrderBy(x => x.Col)} \\[4pt]

SORT (DESC)
  & \texttt{.order\_by(desc(M.col))}
  & \texttt{.orderBy(cb.desc(r.get("col")))}
  & \texttt{.OrderByDescending(x => x.Col)} \\[4pt]

SORT (multi)
  & \texttt{.order\_by(M.a, desc(M.b))}
  & \texttt{.orderBy(cb.asc(...), cb.desc(...))}
  & \texttt{.OrderBy(x=>x.A).ThenByDescending(x=>x.B)} \\[4pt]

LIMIT
  & \texttt{.limit(n)}
  & \texttt{.setMaxResults(n)}
  & \texttt{.Take(n)} \\[4pt]

OFFSET
  & \texttt{.offset(m)}
  & \texttt{.setFirstResult(m)}
  & \texttt{.Skip(m)} \\[4pt]

JOIN (INNER)
  & \texttt{.join(B, A.id==B.a\_id)}
  & \texttt{r.join("bs", JoinType.INNER)}
  & \texttt{join b in Bs on a.Id equals b.AId} \\[4pt]

JOIN (LEFT)
  & \texttt{.outerjoin(B, A.id==B.a\_id)}
  & \texttt{r.join("bs", JoinType.LEFT)}
  & \texttt{join b ... into g from b in g.DefaultIfEmpty()} \\[4pt]

JOIN (CROSS)
  & \texttt{select\_from(A, B)}
  & \texttt{cq.from(A.class); cq.from(B.class)}
  & \texttt{from a in As from b in Bs ...} \\[4pt]

GROUP BY
  & \texttt{.group\_by(M.col)}
  & \texttt{.groupBy(r.get("col"))}
  & \texttt{.GroupBy(x => x.Col)} \\[4pt]

COUNT(*)
  & \texttt{func.count("*")}
  & \texttt{cb.count(root)}
  & \texttt{.Count()} \\[4pt]

SUM
  & \texttt{func.sum(M.col)}
  & \texttt{cb.sum(r.get("col"))}
  & \texttt{.Sum(x => x.Col)} \\[4pt]

AVG
  & \texttt{func.avg(M.col)}
  & \texttt{cb.avg(r.get("col"))}
  & \texttt{.Average(x => x.Col)} \\[4pt]

MIN / MAX
  & \texttt{func.min/max(M.col)}
  & \texttt{cb.min/max(r.get("col"))}
  & \texttt{.Min/.Max(x => x.Col)} \\[4pt]

HAVING
  & \texttt{.having(func.count("*") > n)}
  & \texttt{.having(cb.gt(cb.count(r), n))}
  & \texttt{.Where(g => g.Count() > n)} \\[4pt]

UNION
  & \texttt{q1.union(q2)}
  & \texttt{(JPQL UNION)}
  & \texttt{q1.Union(q2)} \\[4pt]

UNION ALL
  & \texttt{q1.union\_all(q2)}
  & \texttt{(JPQL UNION ALL)}
  & \texttt{q1.Concat(q2)} \\[4pt]

ROW\_NUMBER()
  & \texttt{func.row\_number().over(...)}
  & \texttt{cb.rowNumber().over(...)}
  & \texttt{EF.Functions.RowNumber(...)} \\[4pt]

RANK()
  & \texttt{func.rank().over(...)}
  & \texttt{cb.rank().over(...)}
  & \texttt{EF.Functions.Rank(...)} \\[4pt]

LAG(col, n)
  & \texttt{func.lag(M.col, n).over(...)}
  & \texttt{cb.lag(r.get("col"), n).over(...)}
  & \texttt{EF.Functions.Lag(x.Col, n)...} \\[4pt]

CTE
  & \texttt{subq.cte("name")}
  & \texttt{(Blaze-Persistence CTE)}
  & \texttt{(EF Core 7+ native CTEs)} \\[4pt]

\end{longtable}
\end{landscape}

\begin{thebibliography}{99}

\bibitem{codd1970relational}
E.\,F.\ Codd,
\textit{A Relational Model of Data for Large Shared Data Banks},
Communications of the ACM, 13(6):377--387, 1970.

\bibitem{ansi_sql_2016}
ISO/IEC 9075-1:2016,
\textit{Information technology -- Database languages -- SQL -- Part 1: Framework (SQL/Framework)},
International Organization for Standardization, 2016.

\bibitem{sqlalchemy_docs}
M.\ Bayer et al.,
\textit{SQLAlchemy -- The Python SQL Toolkit and Object Relational Mapper},
\url{https://docs.sqlalchemy.org}, Version 2.0, 2023.

\bibitem{hibernate_docs}
G.\ King, C.\ Bauer, et al.,
\textit{Hibernate ORM Documentation},
\url{https://docs.jboss.org/hibernate/orm/}, Version 6.4, 2023.

\bibitem{ef_core_docs}
Microsoft Corporation,
\textit{Entity Framework Core Documentation},
\url{https://learn.microsoft.com/en-us/ef/core/}, 2023.

\bibitem{linq_docs}
Microsoft Corporation,
\textit{Language Integrated Query (LINQ) -- C\# Programming Guide},
\url{https://learn.microsoft.com/en-us/dotnet/csharp/linq/}, 2023.

\bibitem{java_stream_docs}
Oracle Corporation,
\textit{java.util.stream -- Stream API},
Java SE 17 API Documentation,
\url{https://docs.oracle.com/en/java/documentation/}, 2021.

\bibitem{querydsl_docs}
T.\ Kalliokoski,
\textit{Querydsl -- Unified Queries for Java},
\url{http://querydsl.com}, Version 5.0, 2023.

\bibitem{jooq_docs}
L.\ Waehner,
\textit{jOOQ -- The Best Way to Write SQL in Java},
\url{https://www.jooq.org/doc/}, Version 3.18, 2023.

\bibitem{date_darwen_1997}
C.\,J.\ Date, H.\ Darwen,
\textit{A Guide to the SQL Standard}, 4th Edition,
Addison-Wesley, 1997.

\bibitem{garcia_molina_2009}
H.\ Garcia-Molina, J.\,D.\ Ullman, J.\ Widom,
\textit{Database Systems: The Complete Book}, 2nd Edition,
Pearson Prentice Hall, 2009.

\bibitem{aho_1986}
A.\,V.\ Aho, R.\ Sethi, J.\,D.\ Ullman,
\textit{Compilers: Principles, Techniques, and Tools},
Addison-Wesley, 1986.

\bibitem{wadler1990comprehending}
P.\ Wadler,
\textit{Comprehending Monads},
Proceedings of the 1990 ACM Conference on Lisp and Functional Programming,
pp.\ 61--78, 1990.

\bibitem{buneman1994comprehension}
P.\ Buneman, L.\ Libkin, D.\ Suciu, V.\ Tannen, L.\ Wong,
\textit{Comprehension Syntax},
SIGMOD Record, 23(1):87--96, 1994.

\end{thebibliography}


% ─────────────────────────────────────────────────────────────────────────────
\end{document}
